\documentclass[11pt]{article}
\usepackage[a4paper,margin=25mm,headheight=14pt]{geometry}
\usepackage[T1]{fontenc}
\usepackage{lmodern}
\usepackage{microtype}
\usepackage{amsmath,amssymb,amsthm,mathtools,stmaryrd}
\usepackage{booktabs,longtable,array,tabularx}
\usepackage{enumitem}
\usepackage[dvipsnames]{xcolor}
\usepackage{listings}
\usepackage{xurl}
\usepackage{fancyhdr}
\usepackage{hyperref}
\usepackage[nameinlink,capitalise,noabbrev]{cleveref}

\hypersetup{
  colorlinks=true,
  linkcolor=MidnightBlue,
  citecolor=ForestGreen,
  urlcolor=BrickRed,
  pdftitle={Demand-Driven Relational Certificates for Mathematics in iPRA},
  pdfauthor={OpenAI research synthesis for Vladimir Reshetnikov}
}

\pagestyle{fancy}
\fancyhf{}
\lhead{Demand-Driven Relational Certificates for iPRA}
\rhead{Research report, 15 November 2026}
\cfoot{\thepage}

\newtheorem{theorem}{Theorem}[section]
\newtheorem{lemma}[theorem]{Lemma}
\newtheorem{proposition}[theorem]{Proposition}
\newtheorem{corollary}[theorem]{Corollary}
\newtheorem{definition}[theorem]{Definition}
\newtheorem{principle}[theorem]{Principle}
\newtheorem{example}[theorem]{Example}
\newtheorem{remark}[theorem]{Remark}
\newtheorem{warning}[theorem]{Warning}

\newcommand{\N}{\mathbb N}
\newcommand{\Z}{\mathbb Z}
\newcommand{\Q}{\mathbb Q}
\newcommand{\B}{\{0,1\}}
\newcommand{\ipra}{\mathrm{iPRA}}
\newcommand{\PRAAx}{\texttt{PRAAx}}
\newcommand{\Prf}{\texttt{Prf}}
\newcommand{\NPrf}{\texttt{NPrf}}
\newcommand{\PR}{\mathrm{PR}}
\newcommand{\CFP}{\mathrm{CFP}}
\newcommand{\FUC}{\mathrm{FUC}}
\newcommand{\angles}[1]{\langle #1\rangle}
\newcommand{\abs}[1]{\lvert #1\rvert}
\newcommand{\norm}[1]{\lVert #1\rVert}
\newcommand{\code}[1]{\texttt{\detokenize{#1}}}
\newcommand{\True}{\mathsf{true}}
\newcommand{\False}{\mathsf{false}}
\newcommand{\Viable}{\mathsf{Viable}}
\newcommand{\Local}{\mathsf{Local}}
\newcommand{\Next}{\mathsf{Next}}
\newcommand{\Halt}{\mathsf{Halt}}
\newcommand{\Post}{\mathsf{Post}}
\newcommand{\Ans}{\mathsf{Ans}}
\newcommand{\Run}{\mathsf{Run}}
\newcommand{\Inv}{\mathsf{Inv}}
\newcommand{\Good}{\mathsf{Good}}
\newcommand{\Accept}{\mathsf{Accept}}
\newcommand{\Out}{\mathsf{Out}}
\newcommand{\eps}{\varepsilon}

\lstdefinestyle{leanstyle}{
  language={},
  basicstyle=\ttfamily\small,
  keywordstyle=\color{MidnightBlue}\bfseries,
  commentstyle=\color{ForestGreen},
  stringstyle=\color{BrickRed},
  columns=fullflexible,
  keepspaces=true,
  breaklines=true,
  frame=single,
  framerule=0.3pt,
  rulecolor=\color{gray!55},
  backgroundcolor=\color{gray!4},
  showstringspaces=false,
  tabsize=2
}
\lstset{style=leanstyle}

\title{\textbf{Demand-Driven Relational Certificates for Mathematics in iPRA}\\[0.4em]
\large A compositional extension of finite-observation and certified-finitary frameworks}
\author{Research synthesis prepared for Vladimir Reshetnikov}
\date{15 November 2026}

\begin{document}
\maketitle

\begin{abstract}
The ProveIt iPRA collection contains many mathematically substantial packets, but the same formal obstruction appears repeatedly: a classical proof asks for an exact global choice, an unbounded compactness principle, an unspecified convergence modulus, or a semantic Lean theorem, whereas intuitionistic primitive recursive arithmetic can directly justify only fixed primitive-recursive computations and quantifier-free induction. Existing ProveIt foundation packets already develop bounded Booleanization, finite traces, quantitative names, robust branches, finite nets, and certified finitary presentations. This report consolidates those results and adds a missing compositional layer: \emph{demand-driven relational certificates}. A producer may expose finitely many locally valid answers rather than one discontinuously chosen exact object; a bounded adaptive consumer is compiled together with the producer by a primitive-recursive backward-viability computation. The resulting finite-use composition theorem is proved in detail, including certificate reflection, branch-invariant elimination, quantitative stability, and the precise boundary beyond which an infinite coherent choice would be required.

The framework yields developed human proofs for deterministic trace reflection, relational finite-use composition, guarded arithmetic on represented objects, precision-DAG compilation, ranked metastability, finite-net witness search, branch-invariant algebraic data, and height-bound enumeration. It gives explicit repairs for several recurring difficulties while retaining genuine obstructions: nonzero does not supply apartness, metastability does not supply a final rate, arbitrary Cauchy coefficients do not permit exact Puiseux branch selection, and cardinality bounds do not supply Diophantine search cutoffs. A Lean architecture is specified against the actual \code{ConciseIPRA} API, with precise theorem signatures, substitutions, invariants, and a dependency-ordered completion path. The archive includes a tested exact-arithmetic Python kernel, two Lean source blueprints, a theorem ledger, and a source/difficulty audit. No unproved general Siegel, Faltings, Mordell--Weil, Chebotarev, or prime-distribution theorem is claimed.
\end{abstract}

\tableofcontents
\clearpage

\section{Scope, baseline, and contribution}
\label{sec:scope}

\subsection{The repository situation}

The target theory in the repository is ordinary first-order intuitionistic primitive recursive arithmetic: equality and intuitionistic first-order logic over a language containing a symbol for each fixed primitive-recursive function, the defining equations of those symbols, the successor axioms, and induction for quantifier-free formulas.  The checked \code{ConciseIPRA} library exposes this theory as terms \code{PTm}, symbols \code{PR}, axioms \code{PRAAx}, and derivations \code{Prf PRAAx}.  Its theorem \code{pra_iff_native} transports a completed ordinary iPRA derivation to the native scalar calculus.

The surrounding proof collection is very different in status.  It contains dozens of research packets spanning algebra, analysis, number theory, geometry, probability, and topology.  Their useful mathematical content is considerable, but most are reports, exact-arithmetic experiments, semantic Lean developments, or formalization plans.  They are not interchangeable with inhabitants of \code{Prf PRAAx}.  In particular, the existence of a host-Lean theorem about \code{Nat}, \code{Rat}, polynomials, or Cauchy sequences does not by itself establish an object-theory theorem, even if its program is computable.

Seven foundation packets already identify most of the recurring techniques:
\begin{enumerate}[label=(\roman*)]
\item bounded formulas are compiled to primitive-recursive Boolean tests;
\item bounded loops are verified by quantifier-free invariants;
\item represented objects are used through finite precision queries and explicit moduli;
\item division, normalization, and branch selection are guarded by positive quantitative data;
\item compactness is replaced by finite nets and checked refinements;
\item Diophantine finiteness is separated into a finite verifier and a missing height or covering producer;
\item exact-choice and non-effectivity barriers are recorded rather than hidden by coding.
\end{enumerate}
The most recent baseline is the Certified Finitary Presentation (CFP) packet.  It states a finite-trace reflection theorem and an architecture for object-level Lean, but leaves the generic \code{Prf PRAAx} construction unfinished.  Another packet develops finite relational outputs and a backward dynamic program, but stops short of a general theorem saying how such a multivalued producer composes with an adaptive downstream proof.

\subsection{What is added here}

The main new organizing theorem is a \emph{finite-use relational composition principle}.  Its input is not necessarily a function.  It is a primitive-recursive, finite relation of admissible answers, with local evidence and compatibility checks.  A downstream consumer may ask bounded, adaptive questions: its next query can depend on earlier answers.  Instead of selecting a global exact branch, one forms the finite interaction tree induced by the actual consumer and computes, backward, which states admit a successful continuation.  A forward bounded search then returns the least viable answer at each stage.  All tests are Boolean, all bounds are explicit, and the soundness invariant is quantifier free.

This theorem has four consequences that are useful across the repository.

\begin{enumerate}[label=\arabic*.]
\item \textbf{A finite consumer does not need a global selector.}  Approximate roots, projective directions, local charts, and finite witness families may remain relational as long as the final proof uses only finitely many responses.
\item \textbf{Branch-invariant conclusions can be exported exactly.}  If every admissible transcript gives the same coded output, or gives outputs provably equivalent under the requested observation, the arbitrary least viable transcript may be erased from the final theorem.
\item \textbf{Quantitative stability replaces coherence when appropriate.}  If outputs from all admissible transcripts lie within a supplied error diameter, any transcript gives the required finite observation; no compatible infinite path is needed.
\item \textbf{The obstruction boundary is explicit.}  When a theorem really demands an individually named exact branch at every precision, the finite-use theorem does not manufacture it.  The missing object is exposed as an additional choice, separation, or exact-algebraic certificate.
\end{enumerate}

The framework also sharpens the path to Lean.  It gives a source-level implementation of the generic induction wrapper over the current \code{PRA.qfInd} API, an exact record for relational consumers, and theorem signatures that can be completed without redesigning the mathematics.

\subsection{Claim discipline}

The report uses three status labels.

\begin{description}[style=nextline]
\item[Human theorem.] A complete mathematical proof is supplied, with primitive-recursive bounds and an induction audit.
\item[Semantic executable.] A finite reference implementation has been run on exact data.  This is evidence about the algorithm, not an object-iPRA proof.
\item[Embedded target.] A precise \code{Prf PRAAx} statement and construction path are given.  Unless explicitly stated otherwise, the generic object-language files in this archive have not been compiler-checked against the full repository.
\end{description}

No statement in this report should be read as a proof of a general ineffective theorem merely because the downstream verification is finite.  A certificate checker proves: \emph{if a certificate satisfying these finite equations is supplied, then the conclusion follows}.  Producing the certificate may contain the substantive unresolved mathematics.

\subsection{Road map}

\Cref{sec:logic} fixes the object-theory discipline.  \Cref{sec:diagnosis} classifies the recurring failures.  \Cref{sec:deterministic} develops deterministic finite traces.  \Cref{sec:relational} proves the finite-use relational composition theorem.  \Cref{sec:quantitative,sec:precision} build represented arithmetic and precision calculus.  \Cref{sec:algebraic,sec:compactness,sec:domains} apply the framework to algebraic choice, compactness, geometry, analysis, and Diophantine search.  \Cref{sec:lean} gives the Lean plan, and \Cref{sec:examples} works through exact finite examples.  The final sections record status, limitations, and further questions.

\section{Logical target and compilation discipline}
\label{sec:logic}

\subsection{The object theory}

Fix a finite tuple of numerical parameters $x\in\N^r$.  A \emph{fixed primitive-recursive symbol} means a symbol obtained from zero, successor, projections, composition, and primitive recursion.  The syntax may contain one such symbol for each definition used in a theorem.  It does not contain a single total evaluator that takes an arbitrary program code and runs it without a clock.  This fixed-program convention is essential: a universal unclocked evaluator for all primitive-recursive programs cannot itself be primitive recursive, by diagonalization.

The permitted induction rule has the schematic form
\[
  A(0,x)\;\wedge\;\forall i\bigl(A(i,x)\to A(i+1,x)\bigr)
  \;\Longrightarrow\;\forall i\,A(i,x),
\]
where $A$ is quantifier free.  Parameters $x$ may remain free.  Thus the common phrase ``iPRA cannot reason about a Cauchy name because its contract is universally quantified'' is too coarse.  The global name contract may remain as a hypothesis outside the induction predicate.  At stage $i$, universal elimination instantiates it at the finitely computed query indices required by the current state.  The induction invariant itself records only the resulting Boolean facts.

\subsection{Bounded Booleanization}

A relation is \emph{PR-decidable} when it is represented by a fixed Boolean-valued primitive-recursive term.  Let bounded formulas be generated from PR-decidable atoms using propositional connectives and bounded quantifiers
\[
  \exists j<b(x)\,\varphi(x,j),\qquad
  \forall j<b(x)\,\varphi(x,j).
\]

\begin{theorem}[Bounded Booleanization]
\label{thm:booleanization}
For every fixed bounded formula $\varphi(x)$ there is a Boolean-valued primitive-recursive function $\llbracket\varphi\rrbracket(x)$ such that iPRA proves
\[
  \varphi(x)\;\leftrightarrow\;\llbracket\varphi\rrbracket(x)=1.
\]
Moreover, the construction is compositional and produces a quantifier-free equation suitable for an induction invariant.
\end{theorem}

\begin{proof}
Proceed by structural recursion on $\varphi$.  Atomic formulas already have Boolean tests.  Negation, conjunction, disjunction, and implication are represented by the corresponding truth-table functions on $\{0,1\}$.  For a bounded existential define the prefix fold
\[
 E(x,0)=0,\qquad
 E(x,n+1)=\max(E(x,n),\llbracket\psi\rrbracket(x,n)).
\]
Quantifier-free induction on $n$ proves
\[
 E(x,n)=1\leftrightarrow \exists j<n\,\psi(x,j).
\]
The successor step is a finite case split on the two Boolean values.  Substitute $n=b(x)$.  Bounded universal quantification is compiled either by a minimum/product fold or by negating the existential of the Boolean complement.  All definitions are fixed primitive recursions, and the correctness derivations are intuitionistic because the only case distinctions are on numerical Booleans.
\end{proof}

\begin{corollary}[Least bounded witness]
\label{cor:first-witness}
For every Boolean test $P(x,j)$ and bound $b(x)$ there is a primitive-recursive function $s(x)$, with sentinel $b(x)$, such that iPRA proves
\[
\begin{split}
  &s(x)<b(x)\to P(x,s(x))=1,\\
  &s(x)<b(x)\to \forall j<s(x)\;P(x,j)=0,\\
  &s(x)=b(x)\leftrightarrow\forall j<b(x)\;P(x,j)=0.
\end{split}
\]
Each displayed bounded universal may itself be represented by its Boolean fold.
\end{corollary}

\begin{proof}
Maintain a register initially equal to the sentinel.  At stage $j$, if the register is still the sentinel and $P(x,j)=1$, replace it by $j$; otherwise retain it.  The invariant says that the register is the least successful index in the processed prefix, or the sentinel if none exists.  Booleanize the bounded clauses and apply quantifier-free induction.
\end{proof}

\subsection{Totalization and guards}

Primitive-recursive programs are total.  Mathematical operations such as reciprocal, normalization, or decoding of a nonempty list are partial.  The correct compilation is therefore:
\begin{enumerate}[label=(\alph*)]
\item define a total program with a harmless default outside its domain;
\item define a Boolean guard certifying that the intended domain condition holds;
\item prove the mathematical specification under that guard.
\end{enumerate}
For instance, an integer-coded rational reciprocal may return $0$ when the numerator is $0$.  Every theorem using it carries a positive lower-bound or nonzero integer premise that proves the guard.  Totalization never licenses cancellation without the guard.

\subsection{Fixed precision, then universal introduction}

Suppose a theorem about a represented object has the external form
\[
  H(p)\to\forall k\;\Out(p,k,A(p,k)).
\]
The preferred derivation keeps $p$ and $k$ free, constructs a finite trace specialized to that $k$, and proves $\Out(p,k,A(p,k))$.  Only after the finite proof is complete is $k$ universally introduced.  Induction is never applied to a formula already containing the completed $\forall k$ conclusion.  This order is not cosmetic: it is the mechanism by which all-precision statements remain compatible with quantifier-free induction.

\subsection{Semantic proof versus embedded proof}

A semantic theorem can guide the object proof, but the bridge has four explicit steps:
\begin{enumerate}
\item choose concrete \code{PR} symbols and \code{PTm} terms;
\item formulate every loop invariant as a quantifier-free \code{Fm PTm};
\item construct base and step derivations using defining equations and \code{PRA.qfInd};
\item only then apply \code{pra_iff_native} if a native proof object is wanted.
\end{enumerate}
The final bridge is faithful, but it is not a reflection tactic that turns arbitrary host-Lean propositions into object formulas.

\section{Diagnosis of recurring difficulties}
\label{sec:diagnosis}

The packets exhibit several superficially different failures that reduce to a small number of interface mismatches.  The following classification matters because only some failures are repaired by better coding.

\begin{longtable}{@{}p{0.20\textwidth}p{0.28\textwidth}p{0.24\textwidth}p{0.20\textwidth}@{}}
\toprule
Difficulty & Invalid classical move & Constructive replacement & Residual obligation \\
\midrule
\endhead
Unbounded-precision induction & Put a Cauchy or uniformity quantifier inside the induction invariant & Fix the requested precision; instantiate only the finite query table; quantify afterward & Prove a primitive-recursive query bound and local error inequalities \\
\addlinespace
Exact boundary choice & Decide $x<0$, $x=0$, or $x>0$ for a represented real & Use overlapping branches; exact branch only with a positive gap & Supply a margin or accept a relational output \\
\addlinespace
Nonzero used as invertible & Infer a numerical inverse modulus from $\neg(x=0)$ & Carry an apartness witness $\abs{x}\ge 2^{-m}$ & Prove or assume the quantitative gap \\
\addlinespace
Abstract compactness & Select a convergent subsequence or finite subcover without data & Supply a finite net, coverage modulus, and bounded search & Prove coverage and stability on the domain \\
\addlinespace
Exact root/eigenvector ordering & Globally label branches across collisions & Root clusters, projective data, or finite relational candidates & Isolation/separation if an exact label is required \\
\addlinespace
Metastability upgraded to convergence rate & Treat one stable window as eventual stability & Use the window only locally, or supply a tail/rank certificate & A true final modulus is extra data \\
\addlinespace
Cardinality upgraded to search cutoff & Enumerate until the asserted number of points has appeared & Supply a height bound or a complete-list certificate & Effective height/covering theorem \\
\addlinespace
Ambient Lean theorem called iPRA & Rely on computability-looking definitions in the metatheory & Reify programs and build \code{Prf PRAAx} derivations & Object-level Boolean/rational libraries and invariants \\
\bottomrule
\end{longtable}

\subsection{Logical obstructions versus missing estimates}

It is useful to distinguish three kinds of open edge.

\begin{description}[style=nextline]
\item[Type I: finite formalization debt.]  The mathematics is already a bounded calculation, but no reusable object-language library exists.  Examples include Boolean folds, rational cross-multiplication, finite sums, and trace reflection.  These are difficult engineering tasks but not conceptual obstructions.

\item[Type II: missing quantitative hypotheses.]  A classical statement is true only after silently invoking exact choice or a qualitative condition.  The constructive theorem becomes valid once supplied with an apartness, isolation, modulus, compactness, or chart certificate.  Root selection at a simple separated root is the standard example.

\item[Type III: missing mathematics or genuine nonuniformity.]  No known effective height bound, zero-free constant, or global selector is available, or a proposed selector would imply a principle such as LLPO/LPO or decide a halting-encoded family.  No encoding layer can replace the absent theorem.
\end{description}

The framework is designed to eliminate Type I debt, make Type II inputs explicit and compositional, and expose Type III boundaries in theorem statements.

\subsection{A diagnostic test for a proposed theorem}

Before formalization, write the intended output as a data type and ask four questions.

\begin{enumerate}[label=\arabic*.]
\item Is the output uniquely and continuously determined by the represented input?  If not, can it be replaced by a finite relation, a projective class, a cluster, or a branch-invariant observation?
\item Does every use of division, normalization, sign, or ordering carry a rational margin that a finite checker can verify?
\item Does every compactness or convergence step expose a primitive-recursive finite bound, net, modulus, rank, or certificate?
\item Is the substantive producer known, or has only its finite verifier been described?
\end{enumerate}
A negative answer to the first three suggests a reformulation.  A negative answer to the fourth identifies a genuine unresolved theorem rather than a formalization problem.

\subsection{Why finite relational data are the missing middle layer}

The existing deterministic trace framework expects every upstream component to return one code.  The evidence framework permits finite compatible families but does not explicitly characterize the downstream computations that can safely consume them.  Many packets lie exactly between these interfaces.  A proof of an identity may need one approximate root for five bounded computations but need not name the same root at every future precision.  A projective theorem may need a direction only up to sign.  A symmetric polynomial theorem may use all roots but be invariant under permutation.  For these cases, demanding a global selector introduces a discontinuity that is absent from the mathematical conclusion.

The correct middle layer is therefore not ``choose an infinite branch'' but ``compile the finite interaction required by this consumer.''  The next two sections first recall deterministic traces and then prove the relational extension.

\section{Deterministic finite traces}
\label{sec:deterministic}

\subsection{Trace systems}

A deterministic finite trace system is a tuple of fixed primitive-recursive functions
\[
\begin{array}{lll}
  \Good(u)\in\B, & \mathsf{init}(u), & \mathsf{bound}(u),\\
  \mathsf{step}(u,i,s), & \Inv(u,i,s)\in\B, & \Accept(u,s)\in\B.
\end{array}
\]
All theorem parameters, the requested output precision, and all finite external certificates are paired into $u$.  Define
\[
\Run(u,0)=\mathsf{init}(u),\qquad
\Run(u,i+1)=\mathsf{step}(u,i,\Run(u,i)).
\]
The local obligations are
\begin{align}
 \Good(u)=1&\to \Inv(u,0,\mathsf{init}(u))=1, \tag{D0}\label{eq:D0}\\
 \Good(u)=1\wedge \Inv(u,i,s)=1&\to
   \Inv(u,i+1,\mathsf{step}(u,i,s))=1, \tag{D1}\label{eq:D1}\\
 \Good(u)=1\wedge \Inv(u,\mathsf{bound}(u),s)=1&\to
   \Accept(u,s)=1. \tag{D2}\label{eq:D2}
\end{align}
Each formula is quantifier free after Boolean compilation.

\begin{theorem}[Deterministic trace reflection]
\label{thm:det-trace}
If iPRA proves \eqref{eq:D0}--\eqref{eq:D2}, then it proves
\[
 \forall u\;\bigl(\Good(u)=1\to
 \Accept(u,\Run(u,\mathsf{bound}(u)))=1\bigr).
\]
\end{theorem}

\begin{proof}
Keep $u$ free and apply quantifier-free induction to
\[
 A(i)\;\equiv\;\Good(u)=1\to\Inv(u,i,\Run(u,i))=1.
\]
For $i=0$, rewrite the defining equation for $\Run$ and use \eqref{eq:D0}.  For the successor step, assume $A(i)$ and $\Good(u)=1$.  Modus ponens gives the invariant for $\Run(u,i)$; instantiate \eqref{eq:D1} at $s=\Run(u,i)$; then rewrite $\Run(u,i+1)$ by its defining equation.  Specialize the resulting universal formula at $\mathsf{bound}(u)$ and use \eqref{eq:D2}.  Finally introduce $u$ universally.
\end{proof}

The theorem is deliberately a proof-composition result, not a completeness theorem.  It does not prove \eqref{eq:D0}--\eqref{eq:D2}.  Those are precisely the theorem-specific arithmetic obligations.

\subsection{Trace certificates}

An external producer may return an encoded list $c=(s_0,\ldots,s_B)$.  A fixed Boolean checker $\mathsf{TraceOK}(u,c)$ verifies
\[
\begin{gathered}
B=\mathsf{bound}(u),\qquad s_0=\mathsf{init}(u),\\
s_{i+1}=\mathsf{step}(u,i,s_i)\quad(i<B),\qquad
\Inv(u,i,s_i)=1\quad(i\le B),\\
\Accept(u,s_B)=1.
\end{gathered}
\]
All bounded conjunctions are Boolean folds.

\begin{proposition}[Certificate reflection]
\label{prop:det-cert-reflection}
For a fixed trace language, iPRA proves
\[
 \mathsf{TraceOK}(u,c)=1\to\Accept(u,\mathsf{last}(c))=1,
\]
and, if the transitions are included, also proves
$\mathsf{last}(c)=\Run(u,\mathsf{bound}(u))$.
\end{proposition}

\begin{proof}
Fold elimination yields each checked field.  The acceptance conclusion is immediate.  For equality with the recomputed run, use quantifier-free induction on the list index with invariant $s_i=\Run(u,i)$.  The base is the checked initialization equation.  The step combines the checked transition with congruence and the defining equation of $\Run$.
\end{proof}

Certificate mode separates a possibly expensive search from a small trusted checker.  It is suitable for polynomial identities, interval subdivisions, finite covers, and numerical constant tables.  It never turns an unverified certificate producer into a theorem.

\subsection{Adaptive finite queries to global hypotheses}

Suppose a represented input is given by an evaluator $a(p,n)$ and a global regularity formula
\[
 \mathsf{Regular}(p)\;\equiv\;\forall m\,\forall n\;R(p,m,n,a(p,m),a(p,n)).
\]
The trace may compute its next query from earlier samples.  Let the state contain a table of all queried indices and answers.  At stage $i$, the query term $q(u,i,s)$ is a primitive-recursive function of the current state.  Universal elimination instantiates $\mathsf{Regular}(p)$ at the numerical terms appearing in the state.  The invariant stores the resulting finite Boolean facts.

\begin{lemma}[Adaptive finite instantiation]
\label{lem:adaptive-instantiation}
Let a bounded trace make at most $B(u)$ adaptive queries to a universally quantified PR-decidable contract.  There is a quantifier-free table invariant asserting all contract instances used so far, and iPRA proves that the global contract preserves this invariant at every trace step.
\end{lemma}

\begin{proof}
At stage zero the table is empty.  Assume the table is correct for the processed queries.  The next index is a term in the current numerical state, so universal elimination may instantiate the global contract at that term and at every previously stored term with which it must be compared.  Insert the new Boolean instances into the encoded table.  A bounded fold checks that old entries are unchanged and new entries satisfy the instantiated contract.  This update is primitive recursive, and the invariant is one Boolean equation.  Quantifier-free induction proves it for all stages below $B(u)$.
\end{proof}

The lemma explains why a global all-precision hypothesis is compatible with a finite proof.  The proof consumes only the finite portion demanded by the trace.

\subsection{Sequential composition}

Two deterministic trace systems compose by adding a phase tag.  Phase zero runs the first system.  At its terminal state, the second system is initialized from the first output and from certified metadata such as magnitude bounds or domain guards.  Phase one runs the second system.

\begin{theorem}[Deterministic transducer composition]
\label{thm:det-compose}
If $T_1$ and $T_2$ have iPRA-verified local obligations, and the terminal certificate of $T_1$ implies the input guard of $T_2$, then their sequential composition is a deterministic trace system with an iPRA-verified soundness theorem.  Its bound is the sum of the component bounds plus one switch step.
\end{theorem}

\begin{proof}
Encode the composite state as $(\mathsf{phase},s_1,s_2,m)$, where $m$ is frozen metadata.  The invariant is a Boolean case split.  In phase zero it is the first invariant.  At the switch it contains first acceptance and second input validity.  In phase one it is the second invariant conjoined with preservation of $m$.  Every branch is quantifier free.  The local proof is a finite case analysis on the phase and on whether the current counter is the switch point.  Apply \Cref{thm:det-trace}.
\end{proof}

This theorem handles deterministic proof networks.  It does not yet handle a producer whose mathematically correct output is a finite set of alternatives.  Forcing such a producer into a function is exactly where many discontinuous choices enter.  The next section removes that restriction.

\section{Demand-driven relational certificates}
\label{sec:relational}

\subsection{Bounded evidence machines}

A \emph{bounded evidence machine} describes a finite interaction between a consumer and a relational producer.  Fix an input code $x$ and a resource bound $h$.  The machine has fixed primitive-recursive functions
\[
\begin{array}{ll}
  \mathsf{start}(x), & \Halt(x,r,s)\in\B,\\
  \mathsf{answerBound}(x,r,s), & \Local(x,r,s,a)\in\B,\\
  \Next(x,r,s,a), & \Post(x,r,s)\in\B,
\end{array}
\]
where $r$ is the number of remaining interaction steps and $s$ is the complete finite consumer state, including any history needed to check compatibility.  If $\Halt(x,r,s)=1$, the consumer is finished and $\Post$ is the desired finite conclusion.  Otherwise it accepts an answer $a<\mathsf{answerBound}(x,r,s)$ satisfying $\Local=1$ and moves to $\Next(x,r,s,a)$.

The relation $\Local$ may include several kinds of evidence:
\begin{itemize}
\item residual bounds for an approximate root;
\item compatibility with earlier answers stored in $s$;
\item chart or pivot validity;
\item an apartness margin for a denominator;
\item membership in a finite cover or net cell;
\item a certificate fragment supplied by an external producer.
\end{itemize}
No assumption is made that there is a single answer function valid for all states or all horizons.

\begin{definition}[Backward viability]
\label{def:viability}
Write $A_{r,s}=\mathsf{answerBound}(x,r+1,s)$ and
$s_a=\Next(x,r+1,s,a)$.  Define the Boolean $\Viable(x,r,s)$ by primitive recursion on $r$:
\[
\Viable(x,0,s)=\Halt(x,0,s)\cdot\Post(x,0,s),
\]
and
\[
\begin{aligned}
\Viable(x,r+1,s)=\max\bigl\{&\Halt(x,r+1,s)\cdot\Post(x,r+1,s),\\
 &(1-\Halt(x,r+1,s))
   \max_{a<A_{r,s}}
   \bigl(\Local(x,r+1,s,a)\cdot\Viable(x,r,s_a)\bigr)\bigr\}.
\end{aligned}
\]
Every maximum is a bounded Boolean fold.
\end{definition}

The use of a state rather than a bare level makes the theorem adaptive.  The next query and answer set may depend on all previous replies, provided that history has a bounded numerical code.

\subsection{Semantic correctness of viability}

Call a finite transcript from $(r,s)$ \emph{successful} if it either halts immediately with $\Post=1$, or begins with a locally valid answer and continues successfully from the next state with one fewer remaining step.

\begin{theorem}[Finite-use viability]
\label{thm:viability}
For every fixed bounded evidence machine, iPRA proves, uniformly in $x,r,s$,
\[
 \Viable(x,r,s)=1
 \quad\leftrightarrow\quad
 \text{there is a successful transcript of length at most $r$ from $s$}.
\]
The right side is a bounded formula and may be replaced by its Booleanization.
\end{theorem}

\begin{proof}
Induct on $r$.  At $r=0$, the only permitted transcript has no answer steps, so success is exactly $\Halt(x,0,s)=1$ and $\Post(x,0,s)=1$, the defining equation for viability.

For $r+1$, split on the numerical Boolean $\Halt(x,r+1,s)$.  If it is $1$, success is exactly the terminal postcondition, represented by the first term in the outer maximum.  If it is $0$, a successful transcript must choose an answer $a<A(x,r+1,s)$ for which $\Local=1$, followed by a successful transcript from $\Next(x,r+1,s,a)$ with resource $r$.  By the induction hypothesis, the latter condition is equivalent to the corresponding viability Boolean.  Bounded Booleanization converts the existential over $a$ into the inner maximum.  The converse reconstructs the first answer and appends the inductively supplied tail.  The induction predicate is an equality between Boolean terms, hence quantifier free.
\end{proof}

\subsection{A primitive-recursive transcript selector}

Define $\mathsf{choose}(x,r,s)$ to be the least $a<A(x,r,s)$ satisfying
\[
 \Local(x,r,s,a)=1\quad\text{and}\quad
 \Viable(x,r-1,\Next(x,r,s,a))=1,
\]
with a default when no such $a$ exists.  Define $\mathsf{exec}$ by primitive recursion, halting when the terminal branch is selected and otherwise appending the chosen answer.

\begin{theorem}[Finite-use relational composition]
\label{thm:finite-use-compose}
If
\[
 \Viable(x,h,\mathsf{start}(x))=1,
\]
then the primitive-recursive execution $\mathsf{exec}(x,h)$ returns a finite transcript whose every answer satisfies the advertised local relation and whose terminal state satisfies $\Post=1$.  The theorem is provable in iPRA by quantifier-free induction.  Consequently, a relational producer can be composed with any bounded adaptive consumer once initial viability has been certified.
\end{theorem}

\begin{proof}
The least-witness theorem supplies a chosen answer whenever the nonterminal viability equation is $1$.  Use induction on the number of executed steps, with invariant:
\begin{enumerate}[label=(\roman*)]
\item the encoded transcript length is the number of processed answer steps;
\item every stored edge satisfies $\Local=1$;
\item the current state is obtained by replaying $\Next$ along that transcript;
\item the current state is viable for the remaining resource.
\end{enumerate}
All four clauses are bounded equations or Boolean folds.  Initially the transcript is empty and viability is the hypothesis.  At a terminal state, viability implies $\Post=1$.  At a nonterminal state, unfold viability.  The bounded existential fold is $1$, so the least-witness selector returns a locally valid answer whose successor is viable.  Append it and preserve the invariant.  Because the resource decreases at every answer step, the execution halts within $h$ steps.  The final invariant and the terminal case give the postcondition.
\end{proof}

This theorem is the relational analogue of deterministic transducer composition.  It does not construct an infinite branch, nor does it assert that choices made for two different resource bounds agree.

\subsection{Certificate reflection for interaction trees}

Computing the whole viability table inside a proof can be large.  A certificate may instead contain a finite directed acyclic graph of visited states, one Boolean viability label per node, a selected outgoing edge for each live node, and terminal postcondition checks.

\begin{definition}[Viability certificate]
A code $c$ is valid for $(x,h)$ when a Boolean checker verifies:
\begin{enumerate}[label=(\alph*)]
\item every node carries a resource $r\le h$ and a state code;
\item terminal labels agree with $\Halt\cdot\Post$;
\item every node labeled viable and nonterminal has a listed answer satisfying $\Local$, pointing to a node with resource $r-1$ labeled viable;
\item the root is $(h,\mathsf{start}(x))$ and is labeled viable;
\item replay of the selected edges terminates at a postcondition node.
\end{enumerate}
\end{definition}

\begin{proposition}[Relational certificate reflection]
\label{prop:rel-cert}
For a fixed evidence-machine language, iPRA proves
\[
 \mathsf{RelCertOK}(x,h,c)=1\to
 \mathsf{TranscriptOK}(x,h,\mathsf{extract}(c))=1.
\]
\end{proposition}

\begin{proof}
Every clause is a bounded Boolean fold over the finite certificate.  Extract the root path by following the selected edge.  Induct on the decreasing resource.  The checker supplies local validity and the pointer to the next viable node.  At resource zero or at a marked terminal node, the checker supplies $\Post=1$.  No search beyond the certificate is required.
\end{proof}

This mode lets an untrusted external solver perform dynamic programming, root clustering, or finite cover search while the object theory checks only a bounded graph.

\subsection{Branch-invariant elimination}

Many theorems do not care which successful transcript was used.  Let $\mathsf{output}(s)$ be a code read from a terminal state, and let $E(x,y,z)$ be a PR-decidable observational equivalence.

\begin{definition}[Observational confluence]
A machine is observationally confluent for $E$ when every pair of successful terminal states reachable from the same input has equivalent outputs:
\[
 \mathsf{TerminalReach}(x,h,t_1)=1\wedge
 \mathsf{TerminalReach}(x,h,t_2)=1
 \to E(x,\mathsf{output}(t_1),\mathsf{output}(t_2))=1.
\]
The reachability predicates are bounded and Booleanized.
\end{definition}

\begin{theorem}[Branch-invariant output]
\label{thm:branch-invariant}
Assume initial viability and observational confluence.  Then the output returned by the least-viable execution is $E$-equivalent to the output of every successful transcript.  If $E$ is coded equality, the final output is unique and the internal choice can be erased from the theorem statement.
\end{theorem}

\begin{proof}
\Cref{thm:finite-use-compose} shows that the selected execution reaches a successful terminal state.  Instantiate observational confluence with this terminal state and an arbitrary successful terminal state.  If $E$ is equality, symmetry and transitivity give a unique code.  If $E$ is projective equivalence, permutation equivalence, or equality up to a requested error, the theorem exports exactly that quotient observation.
\end{proof}

Examples include symmetric polynomials of an unordered root list, a projective line independent of vector scaling, a finite set independent of enumeration order, and an interval enclosure independent of the chosen chart.

\subsection{Quantitative elimination without confluence}

Exact confluence is stronger than necessary for approximate conclusions.  Let $d$ be a coded metric and suppose every successful terminal output lies in a set of diameter at most $\delta(x,h)$.

\begin{theorem}[Stable finite-use output]
\label{thm:stable-output}
If initial viability holds and iPRA proves
\[
 \mathsf{TerminalReach}(t_1)=1\wedge\mathsf{TerminalReach}(t_2)=1
 \to d(\mathsf{output}(t_1),\mathsf{output}(t_2))\le\delta,
\]
then the selected output is within $\delta$ of every admissible output.  If a downstream map has a certified Lipschitz bound $L$, all downstream values differ by at most $L\delta$.
\end{theorem}

\begin{proof}
The selected terminal state is successful.  Apply the diameter hypothesis to it and any other terminal state.  The Lipschitz conclusion is the supplied finite inequality.
\end{proof}

This theorem is often the appropriate constructive substitute for a coherent exact branch.  A finite proof at precision $k$ only needs the diameter to be below its error budget.  Choices at different precisions may differ without affecting the universally quantified finite-observation theorem.

\subsection{The exact limitation}

\begin{warning}[No infinite path principle]
\label{warn:no-infinite-path}
The assertions
\[
 \forall h\;\Viable(x,h,\mathsf{start}(x))=1
\]
and ``there exists one primitive-recursive infinite answer sequence compatible at every level'' are not interchangeable.  The finite-use theorem proves the former type of statement, one free horizon at a time.  It does not derive the latter.
\end{warning}

Lexicographically least transcripts at successive horizons need not be prefix-compatible.  Enforcing compatibility across all horizons is an additional theorem, typically supplied by quantitative uniqueness, nested isolating regions, a contraction, or exact algebraic splitting data.  Where such structure exists, it should be stated as a producer theorem; where it does not, a relational or branch-invariant conclusion is the correct endpoint.

\subsection{Kleisli-style composition law}

It is useful to view a deterministic certificate as a singleton relation and a relational producer as a finite nondeterministic computation.  The bounded consumer theorem defines a finite-use composition operation.  Associativity holds up to transcript re-encoding: flattening two nested bounded interactions into one state machine gives the same set of terminal observations.  The proof is a finite induction on the sum of the resource bounds and uses no choice.  Identity is the one-step relation returning its input.  Thus certified finite relations form a practical compositional layer between raw finite certificates and exact mathematical objects.

\section{Quantitative represented objects and guarded operations}
\label{sec:quantitative}

\subsection{Dyadic regular names}

For concreteness, represent a rational by a normalized integer pair and a real by a fixed evaluator $a(p,n)$ returning a rational code.  A convenient regularity convention is
\[
  \abs{a(p,m)-a(p,n)}\le 2^{-m}+2^{-n}.
  \tag{R}\label{eq:regular-name}
\]
All inequalities are implemented by integer cross-multiplication.  A global hypothesis $\mathsf{Regular}_a(p)$ is the universal closure of \eqref{eq:regular-name}.  The semantic real is only explanatory; internal theorems are stated directly in terms of rational samples and error bounds.

Two regular names $a(p,\cdot)$ and $b(q,\cdot)$ are extensionally equal when
\[
  \forall n\;\abs{a(p,n)-b(q,n)}\le 2^{1-n}.
  \tag{E}\label{eq:name-eq}
\]
The constant is chosen with slack so that reflexivity and transitivity are proved by the triangle inequality and a precision shift.  An apartness certificate is a pair $(m,\sigma)$ with $\sigma\in\{-1,+1\}$ such that
\[
  \sigma\,a(p,m)\ge 2^{-m+2}.
  \tag{A}\label{eq:apartness}
\]
Regularity then gives a fixed positive lower bound on the represented value.  The exact numerical constants may be optimized; their role is to leave a rational verification margin.

\begin{lemma}[Finite sample stability]
\label{lem:sample-stability}
Assume \eqref{eq:regular-name}.  If $m,n\ge k+3$, then
\[
 \abs{a(p,m)-a(p,n)}\le 2^{-k-2}.
\]
Consequently, any finite trace whose query planner requests only indices at least $k+3$ may replace one queried sample by another at a cost within the displayed rational budget.
\end{lemma}

\begin{proof}
By regularity, the difference is at most $2^{-m}+2^{-n}\le2\cdot2^{-k-3}=2^{-k-2}$.  This is an integer power-of-two inequality and can be checked by multiplication.
\end{proof}

\subsection{Addition and negation}

Define the output sample for a sum by
\[
  c(k)=a(k+2)+b(k+2).
\]
Then each input contributes at most $2^{-k-2}$ uncertainty relative to a later sample, so the output is a regular name with the desired $2^{-k}$ scale.  Negation is exact on rational codes.

\begin{proposition}[Addition contract]
\label{prop:add-contract}
There is a fixed primitive-recursive evaluator $\mathsf{addName}$ such that iPRA proves: from the regularity of $a$ and $b$, the evaluator is regular, respects extensional equality, and its $k$th sample differs from the sum of any sufficiently precise input samples by the explicit triangle-inequality budget.
\end{proposition}

\begin{proof}
Keep $k$ free.  Instantiate each global regularity formula at the finitely many indices $k+2,m,n$ occurring in the desired inequality.  Expand the rational addition code and apply the finite triangle inequality.  No induction is needed beyond the arithmetic library; universal introduction over $k$ is the last step.
\end{proof}

\subsection{Products require magnitude bounds}

For products,
\[
  ab-a'b'=(a-a')b+a'(b-b').
\]
A finite error contract therefore needs bounds $\abs{a'},\abs{b}\le M$.  These may be supplied externally or derived from one coarse sample and regularity.  For example, if $\abs{a(0)}\le M_0$, then every later sample is at most $M_0+2$ in magnitude under convention \eqref{eq:regular-name}.  Choose a primitive-recursive integer bound $M$ and request input precision $k+c(M)$ so that
\[
  2M\,2^{-(k+c(M))}\le2^{-k-2}.
\]
A simple choice uses a dyadic exponent exceeding $\log_2(4M+4)$, found by bounded doubling.

\begin{proposition}[Product contract]
\label{prop:product-contract}
From regular names and certified rational magnitude bounds, one can primitive-recursively synthesize an input precision and a rational multiplication trace whose iPRA proof yields a regular product name.  The precision transformer is monotone in the output precision and in the magnitude bounds.
\end{proposition}

\begin{proof}
Use the displayed telescoping identity.  Every factor and error is rational.  The precision transformer ensures that each of the two summands is at most half the allocated budget.  Cross-multiply to eliminate division by powers of two.  Monotonicity follows from monotonicity of the doubling search and of the input error bounds.
\end{proof}

\subsection{Guarded reciprocal}

A qualitative premise $\neg(x=0)$ does not supply a usable numerical denominator.  Let a regular name $a$ carry an apartness certificate proving
\[
  \abs{x}\ge\alpha>0
\]
for a rational $\alpha$.  Request a sample $q$ with $\abs{q-x}\le\alpha/2$.  Then $\abs{q}\ge\alpha/2$.  For another approximation $q'$ to $x$,
\[
 \abs{q^{-1}-(q')^{-1}}
 =\frac{\abs{q-q'}}{\abs{q q'}}
 \le \frac{4\abs{q-q'}}{\alpha^2}.
\]

\begin{theorem}[Apartness-guarded inverse]
\label{thm:guarded-inverse}
There is a total primitive-recursive rational operation $\mathsf{inv}_0$ that returns $0$ at a zero numerator and the exact reciprocal otherwise.  From a regular name and a rational apartness witness $\alpha$, iPRA proves that a suitable precision-shifted application of $\mathsf{inv}_0$ is a regular name of the reciprocal and that the zero-default branch is never taken.
\end{theorem}

\begin{proof}
The query planner chooses an index $n(k,\alpha)$ making the input error at most both $\alpha/2$ and $2^{-k}\alpha^2/16$.  Regularity and the apartness witness imply $\abs{a(n)}\ge\alpha/2$, so the rational numerator is nonzero and the guarded operation uses the reciprocal branch.  For two output samples, apply the displayed difference identity and the lower bounds on both denominators.  The chosen precision makes the result at most the regularity budget.  Every premise is a finite rational comparison; the all-precision theorem is obtained by universal introduction over $k$.
\end{proof}

\subsection{Polynomial evaluation}

Let $P(z)=\sum_{j=0}^d c_jz^j$ and suppose $\abs{z},\abs{w}\le M$.  The identity
\[
 z^j-w^j=(z-w)\sum_{t=0}^{j-1}z^tw^{j-1-t}
\]
gives
\[
 \abs{P(z)-P(w)}
 \le \abs{z-w}\sum_{j=1}^d j\abs{c_j}\max(1,M)^{j-1}.
 \tag{P-Lip}\label{eq:poly-lip}
\]
The right side is a finite integer/rational expression.

\begin{lemma}[Polynomial observation contract]
\label{lem:poly-observation}
For a fixed degree bound and coded coefficients, the Lipschitz majorant in \eqref{eq:poly-lip}, the required input precision, and the rational evaluation trace are primitive recursive.  iPRA proves the resulting residual estimate by finite sum and product induction.
\end{lemma}

\begin{proof}
Prove the power-difference identity by induction on $j$, then the polynomial estimate by induction on the coefficient list.  Both induction predicates are quantifier-free identities/inequalities over coded rationals.  A bounded maximum produces $\max(1,M)$, and a finite fold produces the majorant.
\end{proof}

\subsection{Quotient-aware data}

Several mathematical objects should not be represented by normalized exact coordinates.

\begin{itemize}
\item A projective point is a nonzero vector together with scaling equivalence.  Exact normalization is performed only under an apartness witness for a pivot or norm.
\item A root cluster is a finite region with a root-count certificate, not an ordered list of roots.
\item A subspace is represented by a projector, Pl\"ucker coordinates, or a row-space certificate rather than a chosen basis when eigenvalues collide.
\item A finite set is represented modulo permutation and duplicate removal; downstream functions should state permutation invariance.
\end{itemize}
The branch-invariant theorem of \Cref{thm:branch-invariant} then removes internal finite choices at the correct quotient level.

\subsection{Object-theory interface}

A reusable represented-object module should expose three layers:
\begin{enumerate}[label=\arabic*.]
\item \emph{codes and total evaluators}, all fixed \code{PR} symbols;
\item \emph{Boolean finite observations}, such as regularity instances, interval membership, residuals, and apartness;
\item \emph{embedded contracts}, each a \code{Prf PRAAx} theorem consuming precisely the observations used by the algorithm.
\end{enumerate}
It should not expose a metatheoretic real type as an object-language value.  Semantic real numbers are useful for explaining the theorem, but the formal endpoint is always an arithmetic formula about finite samples.

\section{Precision transformers, error graphs, and metastability}
\label{sec:precision}

\subsection{Demand transformers}

A finite numerical operation should advertise not only its output evaluator but also the input observations it demands.  A deterministic contract has:
\[
  \mathsf{need}(x,k),\qquad
  \mathsf{compute}(x,k,\text{queried data}),\qquad
  \mathsf{error}(x,k)\le2^{-k}.
\]
For multiple inputs, $\mathsf{need}$ is a finite tuple.  The transformer must be primitive recursive and monotone after harmless closure
\[
 \mathsf{need}^{\uparrow}(k)=\max_{j\le k}\mathsf{need}(j).
\]
This closure costs only a bounded maximum and simplifies composition.

\begin{theorem}[Precision-contract composition]
\label{thm:precision-compose}
Let $F$ have demand transformer $N_F$ and error contract $E_F$, and let $G$ have transformer $N_G$ and contract $E_G$.  Suppose the magnitude/domain data produced by $F$ suffice to evaluate $N_G$.  Then $G\circ F$ has primitive-recursive demand transformer
\[
  N_{G\circ F}(k)=N_F(N_G(k),\text{metadata for }F),
\]
possibly enlarged by a finite maximum for the metadata queries, and its error proof is the composition of the two finite contracts.
\end{theorem}

\begin{proof}
At output precision $k$, first compute the precision $n=N_G(k)$ required of the intermediate value.  Request from the original input whatever $F$ requires to certify an $n$-precision intermediate value and its metadata.  The $F$ contract supplies exactly the hypotheses consumed by the $G$ contract.  Substitute the certified finite output into the latter proof.  If $G$ is Lipschitz rather than exact on the intermediate representation, allocate its input error according to the displayed Lipschitz factor.  Closure of primitive-recursive functions under composition and bounded maximum gives the transformer; proof substitution gives the object-theory theorem.
\end{proof}

\subsection{Finite error-budget DAGs}

Large proofs are more naturally represented by a finite acyclic dependency graph.  Each node $v$ computes a quantity $z_v$ from predecessors.  A local certificate gives
\[
  \abs{\Delta z_v}\le \eta_v+
  \sum_{w\to v}L_{v,w}\abs{\Delta z_w},
\]
where $\eta_v$ is local rounding/truncation error and $L_{v,w}$ is a certified nonnegative rational sensitivity bound.

Topologically order the graph and define the propagated budget by forward recursion
\[
  E_v=\eta_v+
  \sum_{w\to v}L_{v,w}E_w.
  \tag{B}\label{eq:budget-recurrence}
\]

\begin{theorem}[Error-DAG soundness]
\label{thm:error-dag}
For a fixed bounded graph code, the budgets $E_v$ in \eqref{eq:budget-recurrence} are primitive recursively computable.  If every local certificate is valid and every source error is below its assigned budget, iPRA proves by finite induction that $\abs{\Delta z_v}\le E_v$ for every node, in particular at the output node.
\end{theorem}

\begin{proof}
The state of the induction contains the table of budgets and the Boolean assertion that all processed node errors satisfy their entries.  At the next node, every predecessor has a smaller topological index.  Substitute the induction bounds into the local inequality and use the defining equation for $E_v$.  The finite sum is a primitive-recursive fold.  Booleanize all rational comparisons and apply quantifier-free induction on the number of processed nodes.
\end{proof}

A backward scheduler can synthesize safe local tolerances.  If all path products and the number of paths are bounded by a rational $C$, choosing every $\eta_v\le2^{-k}/(C(N+1))$ suffices.  The coarse bound is often preferable to a numerically optimal but formally expensive optimization.

\subsection{Finite windows and ranked descent}

A convergence rate gives a final cutoff.  Many constructive arguments provide only a stable finite window.  Let $x_n$ be a coded sequence, $g$ a challenge function, and $\eps_k=2^{-k}$.  A window starting at $n$ is stable when
\[
  \forall i,j\in[n,n+g(n)]\quad d(x_i,x_j)\le\eps_k.
\]
This bounded universal is a Boolean fold once $n,k,g$ are fixed.

\begin{definition}[Ranked instability certificate]
For fixed $k$ and $g$, a rank certificate consists of a primitive-recursive rank $\rho(n)\le R(k,g)$ such that whenever the window at $n$ is not stable, the next candidate
\[
  n'=n+g(n)+1
\]
satisfies $\rho(n')<\rho(n)$.
\end{definition}

\begin{theorem}[Rank-to-metastability]
\label{thm:rank-metastability}
Given a ranked instability certificate, the iteration
\[
 n_0=0,
 \qquad n_{t+1}=n_t+g(n_t)+1
\]
finds a stable window by some $t\le R(k,g)$.  The selector and bound are primitive recursive, and the theorem is provable in iPRA.
\end{theorem}

\begin{proof}
Assume all first $R+1$ windows are unstable.  Repeatedly apply the rank-decrease premise to obtain
\[
 \rho(n_{R+1})<\rho(n_R)<\cdots<\rho(n_0)\le R,
\]
a strictly descending chain of $R+2$ natural numbers starting at most $R$, impossible.  Internally, maintain the quantifier-free invariant
\[
  t\le R+1\;\wedge\;
  \bigl(\text{no stable window among the first $t$ candidates}
  \to \rho(n_t)+t\le R\bigr).
\]
The bounded statement about previous windows is Booleanized.  At $t=R+1$, the consequent would say a natural number is at most $-1$ in the equivalent subtraction-free coding, contradiction.  A first-success search returns the least stable $t$.
\end{proof}

This theorem subsumes common energy-decrement and finite-variation arguments.  The rank may be an integer energy bucket, an exceptional-set count, the number of unresolved cells, or a finite descent measure.

\subsection{Why metastability is not a rate}

\begin{warning}[No silent rate extraction]
\label{warn:meta-rate}
A primitive-recursive bound for every challenged stable window does not in general provide a primitive-recursive $N(k)$ such that all $i,j\ge N(k)$ are close.  The challenge function describes one finite future interval, not the entire tail.
\end{warning}

A final Cauchy rate can be recovered only from additional data, for example:
\begin{enumerate}[label=(\alph*)]
\item a monotone tail majorant $T(n)$ with a primitive-recursive search proving $T(N(k))\le2^{-k}$;
\item a summable finite-variation bound with explicit tail modulus;
\item a contraction with a known ratio and initial radius;
\item a well-founded rank whose exhaustion implies permanent stability, not merely one stable window.
\end{enumerate}
The theorem statement must say which mechanism supplies the final rate.

\subsection{Simultaneous stability}

Suppose $m$ finite-window properties have ranks $\rho_j$ and bounds $R_j$.  Conjoin their Boolean stability tests and use the sum rank
\[
  \rho=\sum_{j<m}\rho_j.
\]
Whenever the conjunction fails, at least one component rank decreases, provided the update rule is designed not to increase the others.  More generally, use a lexicographic or weighted rank with explicit finite bounds.

\begin{corollary}[Common stable window]
\label{cor:common-window}
If a common update candidate decreases a bounded aggregate rank whenever any of finitely many Boolean window conditions fails, then a primitive-recursive bounded search returns a window satisfying all conditions simultaneously.
\end{corollary}

This is the correct way to combine finite exceptional-set, variation, and approximation arguments.  It does not imply simultaneous eventual stability beyond the returned window.

\subsection{Precision synthesis as certificate generation}

The error-DAG theorem suggests an automation boundary.  An untrusted compiler may traverse an expression graph, propose magnitude bounds, apartness guards, node tolerances, and input query indices, and emit a certificate.  The trusted object-language checker verifies:
\begin{itemize}
\item graph acyclicity and topological order;
\item every local arithmetic error inequality;
\item every guard and magnitude bound;
\item the budget recurrence;
\item the final output budget.
\end{itemize}
Thus automatic precision synthesis need not be trusted.  Only the small finite checker and its \code{Prf PRAAx} soundness theorem enter the trusted base.

\section{Algebraic branches, root clusters, and exact-choice barriers}
\label{sec:algebraic}

\subsection{Three distinct input regimes}

Statements about polynomial roots must distinguish three representations.

\begin{enumerate}[label=\arabic*.]
\item \textbf{Exact finite coefficients.}  Integer, rational, or exact algebraic codes provide decidable equality, coefficient normalization, and finite splitting data.  Exact factorization and Newton polygon algorithms can be expressed as bounded certificate computations.
\item \textbf{Regular coefficients with separation data.}  Coefficients are Cauchy names, but the input includes a positive discriminant, derivative, or root-separation witness.  A local exact branch can be named because the selected stratum is quantitatively stable.
\item \textbf{Arbitrary Cauchy coefficients.}  Equality and collisions are not decidable.  A global exact branch labeling is generally discontinuous.  The correct uniform outputs are residual approximations, finite root clusters, multisets, or branch-invariant observations.
\end{enumerate}

Conflating these regimes is a frequent source of overstated constructive claims.

\subsection{Simple-root branches}

Let $P(x,y)$ be a polynomial whose coefficients have regular names.  Suppose the input supplies a center $(x_0,y_0)$, a radius box $B$, a residual bound
\[
  \abs{P(x_0,y_0)}\le\eta,
\]
a derivative gap
\[
  \abs{\partial_yP(x_0,y_0)}\ge\gamma>0,
\]
and finite Lipschitz bounds ensuring that the Newton or contraction map preserves $B$ and has contraction ratio $q<1$.

\begin{theorem}[Certified regular branch]
\label{thm:regular-branch}
From the preceding finite certificate there is a primitive-recursive evaluator for a local branch $y(x)$ on a rational sub-box, with an explicit convergence radius and precision transformer.  iPRA proves each requested residual and uniqueness observation by a finite contraction trace.
\end{theorem}

\begin{proof}
At requested precision $k$, choose the number of contraction steps $N$ so that $q^N R\le2^{-k-c}$.  A bounded doubling/search computation finds such $N$ from the rational certificate.  Evaluate $N$ iterates with input coefficient precision determined by the error-DAG of \Cref{thm:error-dag}.  The invariant states that the current iterate stays in $B$, its distance to any fixed point is at most $q^nR$ plus the accumulated coefficient error, and every denominator guard holds.  This invariant is a Boolean conjunction of rational inequalities.  Quantifier-free induction verifies the trace.  The contraction estimate gives the output observation.  Uniqueness follows from
\[
 d(y_1,y_2)=d(Ty_1,Ty_2)\le qd(y_1,y_2),
\]
so $(1-q)d(y_1,y_2)\le0$, an exact rational inequality.
\end{proof}

The theorem does not discover a regular chart from arbitrary singular input.  Chart discovery is a separate relational or exact-algebraic problem.

\subsection{Finite root clusters}

For an integer or rational polynomial $f$, a \emph{root-cluster certificate} at precision $k$ consists of finitely many rational boxes/disks $D_j$, multiplicities $m_j$, and a finite verification trace proving:
\begin{enumerate}[label=(\alph*)]
\item the regions are pairwise separated by a rational gap;
\item each boundary is root-free according to a chosen exact certificate, such as a Sturm count in the real case or an argument/Rouch\'e certificate in a complex rational setting;
\item $D_j$ contains exactly $m_j$ roots counting multiplicity;
\item the sum of multiplicities is $\deg f$;
\item the diameters satisfy the requested bound.
\end{enumerate}
For Cauchy coefficients, replace exact root counts by robust certificates whose margins dominate coefficient uncertainty.

\begin{theorem}[Cluster refinement as a relational producer]
\label{thm:cluster-refine}
Assume a finite cluster certificate at level $k$.  A bounded subdivision procedure produces a finite relation of valid level-$k+1$ refinements.  Any bounded consumer of root data can be composed with this relation by \Cref{thm:finite-use-compose}.  No global ordering of roots is required.
\end{theorem}

\begin{proof}
Subdivide each rational region into a fixed finite family of children.  Enumerate bounded assignments of multiplicities to children and retain those passing the local count, separation, and coverage checks.  The candidate bound is primitive recursive in the number of children and degree.  The local relation verifies that child multiplicities sum to the parent multiplicity and that each child certificate is sound.  A valid exact or robust subdivision theorem establishes initial viability.  The finite-use composition theorem then supplies exactly the refinement transcript demanded by the consumer.
\end{proof}

The theorem factors the problem cleanly: the generic framework verifies and composes refinements; a domain-specific theorem must prove that an admissible refinement exists.

\subsection{Symmetric and quotient outputs}

Let a consumer take approximate roots $z_1,\ldots,z_d$ but produce a polynomial symmetric in them, a factor product, a power sum, or a multiset predicate.  Different cluster refinements may permute or merge labels, yet the final observation is invariant.

\begin{corollary}[Unordered-root elimination]
\label{cor:unordered-root}
If the consumer output is invariant under every permutation and under every relabeling allowed by the cluster compatibility relation, then the least-viable transcript produces the same final output as any admissible root labeling.  The theorem can be stated without a root-ordering function.
\end{corollary}

\begin{proof}
Take $E$ in \Cref{thm:branch-invariant} to be coded equality of the invariant output.  The algebraic symmetry proof supplies observational confluence.
\end{proof}

This is especially useful for Vieta identities, resultants, discriminants, characteristic polynomials, spectral projectors for separated clusters, and products over all roots.

\subsection{The Puiseux boundary}

Consider
\[
  P_a(X,Y)=Y^2-X^2-a
\]
with a nonnegative represented real $a$.  A uniform operation returning a ramification index and individually accessible exact Puiseux coefficients would reveal whether $a=0$: at the operation's own relevant coefficient index, the branch behavior differs by a fixed positive amount between $a=0$ and $a>0$.  A primitive-recursive Cauchy family may encode nonhalting as $a=0$ and halting as $a>0$.  Hence such an exact selector would decide the encoded halting problem.

\begin{warning}[Exact Puiseux coefficients from arbitrary names]
\label{warn:puiseux}
The relational framework does not, and cannot by finite bookkeeping alone, turn arbitrary Cauchy coefficients into a uniformly named exact Puiseux expansion.  Exact algebraic coefficients, a certified stratum, or additional splitting/separation data are required for that conclusion.
\end{warning}

What can be recovered uniformly is weaker but often sufficient:
\begin{itemize}
\item for each finite $(X,Y)$ precision, a finite list of approximate solution germs with residual bounds;
\item cluster multiplicities and branch-invariant polynomial data;
\item exact local series on a supplied simple-root chart;
\item finite truncations under exact algebraic Newton-polygon certificates.
\end{itemize}

\subsection{A worked relational repair for a singular quadratic family}

Fix rational samples $A(n)$ of $a\ge0$ with error at most $2^{-n}$.  At output precision $k$, select a rational grid for coefficients of a truncated polynomial $y(t)=\sum_{j\le N}c_jt^j$.  A candidate is locally valid when interval arithmetic proves
\[
  \norm{y(t)^2-t^2-A(n)}_{\le N}\le2^{-k}
\]
coefficientwise, with enough slack to dominate the coefficient-name error.  Compatibility between levels means that overlapping coefficients differ by the sum of their error radii.  The candidate set is finite, and a certificate producer may show nonemptiness by interval subdivision.

A consumer that only asks for the residual, for the square $y^2$, or for a symmetric union of branches is branch invariant and therefore composes.  A consumer asking for ``the coefficient that distinguishes $a=0$ from $a>0$'' is intentionally not branch invariant and cannot be justified without stronger input.

\subsection{Factorization and degree-uniformity}

For each fixed coefficient evaluator and input degree $d$, complete approximate factorization may be implemented by a bounded descent on a finite coefficient grid.  Degree-uniformity means that $d$ is an ordinary numerical input to one fixed primitive-recursive evaluator; it does not mean evaluating arbitrary encoded programs.  The trace invariant records:
\begin{enumerate}[label=(\roman*)]
\item the current residual polynomial;
\item the product of extracted factors;
\item a coefficientwise telescoping error bound;
\item a strictly decreasing degree/energy measure;
\item all separation and deflation guards used so far.
\end{enumerate}
Relational clusters permit the factor order to remain arbitrary.  The final product identity is permutation invariant, so \Cref{cor:unordered-root} erases the internal choices.

\subsection{Abel--Ruffini and impossibility statements}

Negative algebraic theorems are often easier to internalize because their finite core is exact: finite groups, explicit resolvent identities, degree arguments, and bounded straightedge-and-compass towers.  The framework contribution is not to reprove these theorems automatically but to make their induction discipline explicit.  A certificate contains the finite field/tower/group data; a checker verifies each algebraic identity; the final contradiction is a quantifier-free arithmetic conclusion.  Any transcendental input, such as $\pi$ in squaring the circle, still needs an explicit apartness/transcendence bound rather than a purely qualitative nonalgebraicity assertion.

\section{Finite compactness, witness search, and coherent names}
\label{sec:compactness}

\subsection{Effective compact presentations}

A compact domain should be presented by finite data rather than by a higher-type compactness axiom.

\begin{definition}[Certified net presentation]
A certified net presentation of a domain $X$ consists of:
\begin{enumerate}[label=(\alph*)]
\item a primitive-recursive finite list $N(k)$ of coded points;
\item a rational radius $\delta_k$;
\item a Boolean membership/domain predicate;
\item an iPRA proof that every admissible input point is within $\delta_k$ of some element of $N(k)$;
\item optional refinement and neighbor relations between $N(k)$ and $N(k+1)$.
\end{enumerate}
\end{definition}
Boxes, simplices, finite probability simplices, rational polytopes, sphere charts, and projective charts each need their own coverage proof.  There is no single abstract compactness axiom in the object theory.

\subsection{Finite-net transfer}

Let $F:X\to\Q$ have a certified modulus: $d(x,y)\le\delta\Rightarrow\abs{F(x)-F(y)}\le\omega(\delta)$.  Suppose a finite net satisfies $F(q)\ge -\eta$ at every $q\in N(k)$.

\begin{theorem}[Finite-net transfer]
\label{thm:finite-net}
If $N(k)$ covers $X$ within $\delta_k$, then iPRA proves
\[
  \forall x\in X\qquad F(x)\ge-\eta-\omega(\delta_k).
\]
The conclusion is understood through the represented finite-observation formulation appropriate to $x$.
\end{theorem}

\begin{proof}
For an arbitrary coded/represented point $x$, the coverage theorem supplies a bounded witness $q\in N(k)$ with $d(x,q)\le\delta_k$.  Apply the modulus to obtain $F(x)\ge F(q)-\omega(\delta_k)$, then the checked finite-net inequality.  Internally, the coverage witness is returned by a bounded search or supplied in a certificate.  All arithmetic is rational.  If $x$ is represented, fix the required output precision first and instantiate only the samples needed to verify membership, distance, and the modulus.
\end{proof}

\subsection{Approximate witness search}

Let $R(x,y)$ be a nonnegative residual on a compact represented domain $Y$, with modulus $\omega$.  Evaluate $R$ on a finite net and choose a least minimizer.

\begin{theorem}[Finite approximate witness]
\label{thm:approx-witness}
Suppose there exists $y_*$ with $R(x,y_*)\le\eta$, and $N(k)$ covers $Y$ within $\delta_k$.  The finite argmin $q_k$ satisfies
\[
  R(x,q_k)\le\eta+\omega(\delta_k)+\text{evaluation error}.
\]
The selector and proof are primitive recursive.
\end{theorem}

\begin{proof}
Choose a net point $q_*$ within $\delta_k$ of $y_*$.  The modulus bounds $R(x,q_*)$ by $\eta+\omega(\delta_k)$.  Approximate evaluation adds the certified finite error.  Since $q_k$ is a checked finite argmin, its evaluated residual is no larger than that of $q_*$.  Transfer the evaluation error once more if needed.  The argmin invariant stores the current least index and the Boolean fold asserting minimality over the processed prefix.
\end{proof}

The theorem does not require uniqueness.  It gives a finite relational output if several net points tie or if approximate comparisons overlap.

\subsection{Nested isolating regions}

A sequence of finite regions can define one coherent represented object when genuine quantitative uniqueness is available.

\begin{definition}[Nested-ball certificate]
A nested-ball certificate consists of centers $c_k$, radii $r_k$, and Boolean proofs that
\[
  B(c_{k+1},r_{k+1})\subseteq B(c_k,r_k),
  \qquad r_k\le2^{-k-c},
\]
and that the target relation has at least one solution in every ball.  A uniqueness modulus asserts that any two admissible solutions in $B(c_k,r_k)$ are within a prescribed diameter.
\end{definition}

\begin{theorem}[Nested balls produce a regular name]
\label{thm:nested-name}
The centers of a nested-ball certificate form a regular name after a fixed precision shift.  If the uniqueness modulus tends to zero with a primitive-recursive rate, every admissible solution is extensionally equal to that name.
\end{theorem}

\begin{proof}
For $m\ge n$, nesting gives $c_m\in B(c_n,r_n)$ up to the rational inclusion convention, hence $d(c_m,c_n)\le r_n+r_m\le2^{-n-c}+2^{-m-c}$.  Choose $c$ so this is below the regularity budget.  For an arbitrary admissible solution, place it in the same isolating region and apply the uniqueness diameter.  All statements at a fixed pair $(m,n)$ are finite; universal introduction follows after the rational proof.
\end{proof}

Nestedness is the missing structure that upgrades separate finite viability at every precision to a coherent infinite name.  It must be proved, not inferred from nonempty candidate sets alone.

\subsection{Strong convexity instance}

If a rationally presented function $F$ satisfies a quantitative strong-convexity inequality
\[
 F\left(\frac{x+y}{2}\right)
 \le\frac{F(x)+F(y)}2-\mu d(x,y)^2
\]
with $\mu>0$, then approximate minimizers are quantitatively close.  If $F(x),F(y)\le I+\eta$, where $I$ is the infimum, then
\[
 \mu d(x,y)^2\le\eta.
\]
Thus finite-net argmin points at successively smaller residuals have a computable diameter and can be turned into nested or regular names.

\begin{corollary}[Coherent minimizer from strong convexity]
\label{cor:strong-minimizer}
An effective compact presentation, a modulus for $F$, a positive strong-convexity constant, and a primitive-recursive error schedule yield a regular name for the unique minimizer.  Without quantitative uniqueness, the framework yields only finite approximate minimizer sets or values of the infimum.
\end{corollary}

\subsection{Finite covers and local charts}

A global theorem may be proved by finitely many charts.  A cover certificate contains charts $U_j$, local guards $G_j$, and a coverage fold asserting every net cell lies in at least one robustly enlarged chart.  Exact chart selection at a boundary is unnecessary: the relational producer retains all charts whose robust guard passes, and the bounded consumer chooses a viable one.  If local conclusions agree on overlaps, \Cref{thm:branch-invariant} exports a chart-independent result.

This pattern applies to sphere atlases, projective geometry, matrix groups, contour subdivisions, and piecewise analytic estimates.  It is safer than a hard ``first chart containing $x$'' rule, which can be discontinuous at chart boundaries.

\subsection{Limits of finite compactness}

Finite nets do not prove a missing modulus, and a nonempty net-level candidate set at every precision does not by itself produce a coherent point.  A valid global conclusion needs at least one of:
\begin{itemize}
\item stability of the desired observation across all candidates;
\item a quantitative uniqueness/isolation modulus;
\item certified nested refinement;
\item an exact finite algebraic representation;
\item a theorem that the particular bounded consumer uses only one finite level.
\end{itemize}
These alternatives correspond exactly to the branch-invariant, stable-output, nested-name, exact-code, and finite-use theorems developed above.

\section{Domain libraries and resolution of reported obstacles}
\label{sec:domains}

\subsection{Geometry and linear algebra}

Coordinate geometry over rationals is naturally finitary.  Incidence, determinants, cross products, and polynomial identities are exact.  The constructive difficulty arises when an algebraic vector must be interpreted as an actual point, line, unit axis, or normalized eigenvector.

The recommended interface is:
\begin{enumerate}[label=\arabic*.]
\item prove the universal polynomial identity without nondegeneracy assumptions;
\item represent geometric objects projectively or by subspaces/projectors;
\item attach explicit apartness to every claimed nonzero coordinate or determinant;
\item use robust overlapping pivots to choose a chart;
\item state exact normalized outputs only under a positive gap.
\end{enumerate}

For an orthogonal rotation matrix, for example, one may always output approximate/projective axis data or a quaternion chart family.  A unique normalized axis with a prescribed sign is requested only when a quantitative spectral or coordinate gap distinguishes it.  At the identity and half-turn boundaries, a relational family is the honest uniform output.  An exact selector over arbitrary represented rotations can encode weak sign choice and therefore cannot be justified merely by a clever normalization formula.

\subsection{Analysis}

Derivatives, integrals, power series, and contour arguments can be compiled when every limiting step has an explicit finite remainder.

\begin{itemize}
\item Riemann integration uses a supplied modulus of continuity, a rational mesh, and upper/lower finite sums.
\item Differentiation uses a finite difference estimate with a quantified remainder and an apartness guard when dividing by an increment.
\item Power series use coefficient majorants and a finite tail bound.
\item Cauchy and residue theorems use rational polygonal contours, denominator separation, finite cell estimates, and a certified winding/root count.
\item Euler--Maclaurin, Stirling, Gamma, and zeta formulas require explicit truncation schedules and rational bounds for every analytic input.
\end{itemize}

The relational layer is useful when the proof admits several contours, saddle regions, or local branches.  A certificate may retain all robust choices; a bounded consumer selects a viable one.  It does not supply a missing zero-free region or an ineffective constant.

\subsection{Analytic number theory}

A reusable analytic-number-theory certificate should separate four layers:
\begin{enumerate}[label=(\alph*)]
\item exact finite arithmetic identities (Dirichlet convolution, divisor sums, Euler factors);
\item finite truncation and tail estimates;
\item zero-free/exceptional-zero or contour certificates;
\item the final finite numerical inequality.
\end{enumerate}
The first and fourth layers are well suited to generic Boolean/rational checkers.  The second is suited when an explicit majorant is known.  The third is often the substantive theorem.  A proof packet must not replace it by a field in a record and then call the record's type a proof.

For prime-distribution results, the framework can verify a supplied zero-free region, density estimate, or explicit-formula error table and compose it with the finite arithmetic downstream.  It does not establish general Chebotarev, Siegel--Walfisz, Bombieri--Vinogradov, or an ineffective exceptional constant on its own.

\subsection{Probability}

Finite rational probability laws are ideal inputs.  Expectations, moments, convolutions, and characteristic functions at rational mesh points are finite sums.  A central-limit proof can be organized as:
\begin{enumerate}[label=\arabic*.]
\item normalize using a positive variance certificate;
\item compute a finite characteristic-function or smoothing trace;
\item use explicit Taylor remainders and moment bounds;
\item evaluate a finite integral/grid estimate;
\item combine errors in a precision DAG.
\end{enumerate}
For nonidentically distributed or martingale variants, the state records the accumulated variance, moment budget, and conditional finite tables.  Any conditioning over an infinite probability space must be replaced by a finite law presentation or an explicit approximation contract.

\subsection{Topology and combinatorics}

Sperner, parity, winding on finite meshes, and discrete Jordan arguments have quantifier-free invariants.  The main obligation is to expose the mesh and separation hypotheses that make the discrete result imply the represented geometric statement.  Exact boundary classification is replaced by staggered grids or overlapping cells.  A finite-use chart consumer is particularly natural in cubical and simplicial subdivisions.

Pure combinatorial identities are usually Type I formalization debt: finite lists, recurrences, determinants, and generating-function coefficients can be handled by Boolean folds and polynomial certificates.  Asymptotic claims still need explicit remainder schedules.

\subsection{Diophantine geometry: height before enumeration}

Let $P(x,z)$ be a fixed PR-decidable predicate for a coded rational or integral point and let $H(z)$ be its height.  A bound $B(x)$ gives a finite enumerator
\[
 L(x)=\{z:H(z)\le B(x)\ \wedge\ P(x,z)=1\}.
\]

\begin{theorem}[Height bound and complete list]
\label{thm:height-list}
For a fixed coded point predicate and height:
\begin{enumerate}[label=(\roman*)]
\item a primitive-recursive height bound yields a primitive-recursive complete finite list by bounded filtering;
\item a primitive-recursive complete finite list yields a primitive-recursive height bound by taking the maximum listed height (with a default for the empty list).
\end{enumerate}
Both implications are provable in iPRA.
\end{theorem}

\begin{proof}
For (i), enumerate the finite box of normalized codes of height at most $B(x)$, test $P$, and retain successful codes.  The completeness proof substitutes the assumed bound for an arbitrary point and shows its normalized code occurs in the box.  For (ii), compute the maximum height in the list.  Completeness says every point equals one of the listed codes, hence its height is at most the maximum.  Finite-list membership, filtering, and maximum correctness are proved by quantifier-free induction.
\end{proof}

A numerical bound on the \emph{number} of points is not enough.  Points may occur at arbitrarily large heights, and seeing the claimed number during a search does not certify that no further point exists unless the list theorem itself supplies completeness.

\subsection{Runge, descent, and power-gap certificates}

Several explicit families in the repository already fit the framework.  A polynomial identity such as
\[
 D^m f(X)=A(X)^m+R(X),
 \qquad \deg R<(m-1)\deg A,
\]
can imply a power gap for sufficiently large $\abs{X}$.  The certificate checker verifies the coefficient identity, degree inequalities, and derived height formula.  Bounded enumeration then lists all solutions.  The finite checker is generic; construction of the identity is a rational polynomial algorithm in Runge-type cases but not a solution of general Siegel.

Descent certificates similarly contain a measure that strictly decreases, local solubility/cover data, and terminal cases.  If a finite covering list is supplied, relational composition can select viable local covers without globally naming representatives.  Producing a uniform covering list or Mordell--Weil generators for every curve remains a separate theorem.

\subsection{Faltings, Siegel, and Mordell--Weil}

The framework improves the statements but not the missing global theorems.

\begin{itemize}
\item For Faltings-type finiteness, it proves the equivalence between a positive height bound and a complete bounded list for the explicit coded predicate.  A cardinality theorem alone does not yield either.
\item For Siegel, it packages Runge and other effective subcases as height-certificate producers, followed by a common verifier.  It does not produce the general height bound.
\item For Mordell--Weil, it verifies descent relative to supplied covers and combines finite local choices.  It does not uniformly compute generators absent an effective covering theorem.
\end{itemize}

These are not failures of the framework.  They are precise factorizations of the remaining mathematics.

\subsection{Resolution matrix}

\begin{longtable}{@{}p{0.25\textwidth}p{0.42\textwidth}p{0.25\textwidth}@{}}
\toprule
Reported difficulty & Resolution supplied here & Status \\
\midrule
\endhead
All-precision hypothesis inside induction & Adaptive finite query table; fix precision, then quantify & Resolved generically \\
\addlinespace
Multivalued approximate witnesses & Backward viability and finite-use relational composition & Resolved for bounded consumers \\
\addlinespace
Branch labels/orderings & Quotient data, confluence, or diameter stability & Resolved when conclusion is invariant/stable \\
\addlinespace
Division from nonzero & Apartness-guarded total inverse & Requires quantitative gap \\
\addlinespace
Compactness/minimizers & Certified nets, argmin, nested isolation & Requires coverage and uniqueness/modulus \\
\addlinespace
Metastability used as a rate & Ranked window theorem plus explicit prohibition of upgrade & Local use resolved; final rate remains extra \\
\addlinespace
Puiseux from arbitrary Cauchy coefficients & Finite residual clusters; exact series only on certified strata & Exact global selector obstructed \\
\addlinespace
Exact rotation axis/sign & Projective/relational output; exact under gap & Unrestricted selector obstructed \\
\addlinespace
General Siegel/Faltings enumeration & Height/list verifier and certificate interface & Height producer unresolved \\
\addlinespace
Semantic Lean proof called iPRA & Explicit \code{Prf PRAAx} module graph and trace wrapper & Engineering milestone, not completed here \\
\bottomrule
\end{longtable}

\section{Lean formalization architecture}
\label{sec:lean}

\subsection{Target the checked repository interfaces}

The implementation should extend the current \code{LeanProofs.ConciseIPRA} library rather than define another arithmetic logic.  The central declarations are:

\begin{longtable}{@{}p{0.37\textwidth}p{0.55\textwidth}@{}}
\toprule
Repository declaration & Use in this framework \\
\midrule
\endhead
\code{PR k} & Fixed primitive-recursive symbols of arity $k$ \\
\code{PTm n} & Arithmetic terms with $n$ free variables \\
\code{Fm PTm n} & First-order formulas over those terms \\
\code{PRAAx} & Defining equations, successor axioms, and quantifier-free induction \\
\code{Prf PRAAx Gamma phi} & The required object-level endpoint \\
\code{PRA.precZero}, \code{PRA.precSucc} & Rewriting a concrete primitive-recursive runner \\
\code{PRA.qfInd} & Quantifier-free invariant induction \\
\code{Prf.subst}, \code{Prf.weaken}, \code{Prf.trans} & Structural proof composition \\
\code{Prf.allE}, \code{Prf.allI} & Specialization at a computed bound and final universal introduction \\
\code{pra_iff_native} & Export only after a \code{PRAAx} derivation exists \\
\bottomrule
\end{longtable}

The ordinary iPRA language is preferable for the first implementation because arithmetic formulas and finite invariants are readable.  Native/public certificate compression is a later optimization.

\subsection{Concrete primitive-recursive trace symbols}

A runner with one parameter has the repository arities
\begin{lstlisting}
structure PRTraceSystem where
  good   : PR 1
  init   : PR 1
  bound  : PR 1
  step   : PR 3   -- predecessor, previous value, parameter
  inv    : PR 3   -- counter, state, parameter
  accept : PR 2   -- state, parameter

def PRTraceSystem.run (T : PRTraceSystem) : PR 2 :=
  .prec T.init T.step
\end{lstlisting}
The order is dictated by \code{PR.prec}: the recursion argument comes first, followed by the ordinary parameters.  Larger inputs and states are paired into one code.

The archive's \code{lean/ObjectTraceReflection.lean} defines term constructors \code{app1}, \code{app2}, and \code{app3}, then the terms
\[
 \Good_T(u),\quad \Run_T(i,u),\quad
 \Inv_T(i,\Run_T(i,u),u),\quad
 \Accept_T(\Run_T(\mathsf{bound}_T(u),u),u).
\]
The loop formula has two free variables, with de Bruijn variable $0$ as the induction counter and variable $1$ as the parameter:
\[
 \phi_T(i,u)=\bigl(\Good_T(u)=1\bigr)\to
 \bigl(\Inv_T(i,\Run_T(i,u),u)=1\bigr).
\]
It is quantifier free by construction.

\subsection{The generic object-level wrapper}

The theorem-specific file should supply exactly three pieces of proof data after all defining-equation rewriting:

\begin{lstlisting}
structure TraceObligations (T : PRTraceSystem) where
  base : [] |-[PRAAx] (tracePhi T).subst (inst PTm.zero)
  step : [tracePhi T] |-[PRAAx]
    (tracePhi T).subst (repl (PTm.succ (Syn.var 0)))
  finish : [phiAtBound T] |-[PRAAx] traceGoal T
\end{lstlisting}

Here \code{phiAtBound T} substitutes \code{T.bound u} for the counter and \code{traceGoal T} is the one-free-variable formula
\[
 \Good_T(u)=1\to
 \Accept_T(\Run_T(\mathsf{bound}_T(u),u),u)=1.
\]
The generic wrapper is mathematically:

\begin{lstlisting}
theorem traceReflection (O : TraceObligations T) :
    [] |-[PRAAx] Fm.all (traceGoal T) := by
  have hall : [] |-[PRAAx] Fm.all (tracePhi T) :=
    PRA.qfInd (tracePhi T) (tracePhi_qf T) O.base O.step
  have hbound : [] |-[PRAAx] phiAtBound T :=
    Prf.allE (boundTerm T) hall
  have hgoal : [] |-[PRAAx] traceGoal T :=
    O.finish.trans (fun psi hpsi => by
      simpa using hbound)
  exact Prf.allI hgoal
\end{lstlisting}

The precise syntax of the final \code{simpa} may require a one-line lemma about singleton membership or substitution normalization, but no new mathematical idea is missing.  The base and step proofs remain theorem specific because they contain the real invariant arithmetic.  The generic wrapper closes the CFP packet's first formalization milestone: the use of \code{PRA.qfInd}, specialization at the computed bound, cut with final soundness, and universal closure are uniform.

\subsection{Rewriting the runner}

For a concrete $T$, the following object equations are derived from the repository library:
\begin{align*}
 \Run_T(0,u)&=\mathsf{init}_T(u),\\
 \Run_T(i+1,u)&=\mathsf{step}_T(i,\Run_T(i,u),u).
\end{align*}
In Lean, instantiate \code{PRA.precZero} and \code{PRA.precSucc} with \code{T.init} and \code{T.step}.  Small helper lemmas normalize the finite argument functions passed to \code{PTm.app}.  The theorem-specific \code{base} proof uses the first equation; \code{step} uses the induction hypothesis, local preservation, the second equation, and equality substitution.

It is useful to expose a second record with raw local obligations (D0)--(D2) and prove a constructor that fills \code{TraceObligations}.  That constructor is more algebraically involved because it must open implications and perform the run rewrites, but it is still generic.  The suggested staging is:
\begin{enumerate}[label=\arabic*.]
\item compile the wrapper above, taking already rewritten \code{base}/\code{step}/\code{finish};
\item implement run-zero and run-successor lemmas;
\item implement implication plumbing from raw local formulas;
\item add a syntax-generating record that constructs those formulas automatically.
\end{enumerate}

\subsection{Boolean folds}

The next module should define fixed symbols and object proofs for:
\begin{itemize}
\item Boolean not/and/or/implication;
\item \code{allBelow}, \code{anyBelow}, and first successful search;
\item finite list fold, lookup with default, length, append, and bounded table update;
\item finite sum, product, maximum, and least-index argmin;
\item pairing and bounded sequence codes.
\end{itemize}

A generic fold theorem has the form
\[
 \mathsf{allBelow}(p,n)=1
 \leftrightarrow \forall i<n\;p(i)=1,
\]
where the right side is used in human statements and the left side in induction invariants.  The internal proof can avoid a literal bounded quantifier entirely by using elimination/intro lemmas for the fold.  For example:
\begin{align*}
 \mathsf{allBelow}(p,n+1)=1&\to\mathsf{allBelow}(p,n)=1,\\
 \mathsf{allBelow}(p,n+1)=1&\to p(n)=1,\\
 \mathsf{allBelow}(p,n)=1\wedge p(n)=1&\to
   \mathsf{allBelow}(p,n+1)=1.
\end{align*}
These equations are sufficient for most object proofs and keep formulas quantifier free.

\subsection{Semantic relational core}

The standalone file \code{lean/FiniteUseSemantic.lean} contains the host-level design:

\begin{lstlisting}
structure EvidenceMachine where
  answerBound : Nat -> Nat -> Nat
  halt        : Nat -> Nat -> Bool
  local       : Nat -> Nat -> Nat -> Bool
  next        : Nat -> Nat -> Nat -> Nat
  post        : Nat -> Nat -> Bool

def viable (M : EvidenceMachine) : Nat -> Nat -> Bool
  | 0, s => M.halt 0 s && M.post 0 s
  | r+1, s =>
      if M.halt (r+1) s then M.post (r+1) s
      else anyBelow (M.answerBound (r+1) s)
        (fun a => M.local (r+1) s a && viable M r (M.next (r+1) s a))
\end{lstlisting}

The file defines proposition-valued targets for viability correctness, selector soundness, transcript replay, confluence, and diameter stability.  It intentionally does not pretend that arbitrary host functions are primitive recursive.  The object module will replace every field by a fixed \code{PR} symbol and every target by a \code{Prf PRAAx} theorem.

\subsection{Object-level relational record}

A first object-language implementation can pair $(r,s,x)$ into one state and use a deterministic trace that fills a finite viability table.  This reduces relational composition to \Cref{thm:det-trace}; no new induction rule is needed.  The record should contain:

\begin{lstlisting}
structure PRRelMachine where
  start       : PR 1
  answerBound : PR 3
  halt        : PR 3
  local       : PR 4
  next        : PR 4
  post        : PR 3
  stateBound  : PR 2
\end{lstlisting}

The additional \code{stateBound} bounds all states encountered at a fixed input/horizon and permits a rectangular dynamic-programming table.  A more compact sparse certificate mode can be added later.

The proof stages are:
\begin{enumerate}[label=\arabic*.]
\item define a table row by bounded reverse recursion;
\item prove the row equation from Boolean-fold correctness;
\item use quantifier-free induction on remaining depth to prove semantic viability;
\item define least viable answer and prove its local/successor properties;
\item define forward execution and apply the generic trace wrapper;
\item prove confluence or diameter stability as a separate bounded theorem;
\item export only the invariant/stable observation.
\end{enumerate}

\subsection{Quantitative libraries}

The recommended module graph is:

\begin{center}
\begin{tabular}{@{}l@{}}
\code{FUC/Codes.lean}\\
$\downarrow$\\
\code{FUC/BooleanFold.lean}\quad \code{FUC/Integer.lean}\quad \code{FUC/Rational.lean}\\
$\downarrow$\\
\code{FUC/TraceReflection.lean}\\
$\downarrow$\\
\code{FUC/Relational.lean}\quad \code{FUC/ErrorDAG.lean}\\
$\downarrow$\\
\code{FUC/RegularNames.lean}\quad \code{FUC/FiniteNet.lean}\\
$\downarrow$\\
\code{FUC/RootCluster.lean}\quad \code{FUC/Projective.lean}\quad
\code{FUC/HeightSearch.lean}\\
$\downarrow$\\
Theorem-specific packet adapters\\
$\downarrow$\\
\code{pra_iff_native} and optional compact public certificates
\end{tabular}
\end{center}

The rational layer should favor direct integer identities over a large abstract ordered-field hierarchy.  The necessary endpoints are cross-multiplied equality/order, finite sums/products, absolute value, powers, guarded inverse, intervals, and polynomial evaluation.  Semantic Mathlib lemmas may validate calculations, but the exported theorem must be about coded natural numbers.

\subsection{An end-to-end proof record}

A theorem packet should expose machine-readable status fields:
\begin{lstlisting}
structure PacketStatus where
  fixedEvaluators      : Bool
  finiteBounds         : Bool
  certificateChecker   : Bool
  humanSoundnessProof  : Bool
  semanticTestsRun     : Bool
  objectDerivation     : Bool
  nativeExport         : Bool
  axiomAuditRun        : Bool
  unresolvedProducer   : Option String
\end{lstlisting}
A record value is metadata, not a proof.  Its purpose is to prevent a compiled specification or a passing finite test from being reported as a completed object theorem.

\subsection{Milestones and audits}

The dependency-ordered plan is:
\begin{enumerate}[label=\arabic*.]
\item compile the generic rewritten trace wrapper in \code{ObjectTraceReflection.lean};
\item add run-zero/run-successor and raw-obligation constructors;
\item compile Boolean folds and least bounded search as \code{Prf PRAAx};
\item prove one recurrence and one robust two-branch example end to end;
\item compile relational viability on a bounded rectangular state space;
\item formalize regular-name addition, product, and apartness-guarded inverse;
\item formalize error-DAG and finite-net argmin;
\item port one substantial packet, preferably a separated Euler--Rodrigues theorem or a finite-law probability theorem;
\item add root clusters and a height/list adapter;
\item only after correctness, optimize proof/certificate size or target the native public calculi.
\end{enumerate}

Every exported endpoint should be checked by:
\begin{itemize}
\item a source scan for \code{sorry}, \code{admit}, and new line-leading \code{axiom};
\item \texttt{\#print axioms} on all public theorems;
\item a clean build against the pinned repository toolchain;
\item independent exact finite regression tests;
\item a theorem ledger distinguishing host semantics, object derivation, and native export.
\end{itemize}

\subsection{Status of the supplied Lean files}

The two Lean files in this archive are formalization blueprints.  \code{FiniteUseSemantic.lean} gives executable-looking host definitions and explicit goal propositions.  \code{ObjectTraceReflection.lean} is written against the inspected \code{ConciseIPRA} API and includes the generic wrapper design.  The artifact environment did not contain Lean or the repository dependency closure, so the files are not claimed compiler-checked.  The report and \code{lean/FORMALIZATION_PLAN.md} identify the exact likely elaboration points and the shortest path to a checked implementation.

\section{Worked finite examples}
\label{sec:examples}

The examples in this section are intentionally small.  Their purpose is to show the shape of a certificate and its invariant, not to replace the larger domain theorems.

\subsection{A relational square-root approximation}

Let $a=2$ and let level $n$ use dyadic candidates
\[
 C_n=\left\{\frac{j}{2^n}:0\le j\le 2^{n+2},\quad
 \abs{\left(\frac{j}{2^n}\right)^2-2}\le 4\cdot2^{-n}\right\}.
\]
Compatibility between $q\in C_n$ and $r\in C_{n+1}$ is
\[
 \abs{q-r}\le 2^{1-n}.
\]
All tests use integer arithmetic after multiplying by $2^{2n}$.

\begin{proposition}
For every finite horizon $h$, the levels $0,\ldots,h$ admit a compatible chain.
\end{proposition}

\begin{proof}
Let $j_n$ be the least integer with $j_n^2\ge2^{2n+1}$ and put $q_n=(j_n-1)/2^n$.  The least-crossing definition is a bounded search because $(2^{n+2})^2$ exceeds the target.  We have
\[
 q_n^2<2\le(q_n+2^{-n})^2,
\]
so
\[
 0<2-q_n^2\le 2q_n2^{-n}+2^{-2n}\le4\cdot2^{-n}
\]
using $q_n<2$.  Thus $q_n\in C_n$.  The intervals
$[q_n,q_n+2^{-n}]$ and $[q_{n+1},q_{n+1}+2^{-(n+1)}]$ both contain $\sqrt2$ in the explanatory metatheory; internally, the same conclusion follows by comparing the crossing integers.  Hence $\abs{q_n-q_{n+1}}\le2^{-n}+2^{-(n+1)}<2^{1-n}$.
\end{proof}

A relational certificate lists the finite candidate sets and the compatible edges.  Backward viability marks candidates that reach level $h$.  Forward selection returns the least viable chain.  A consumer needing only a $2^{-k}$ approximation asks for one sufficiently fine level; it does not need a global exact square-root selector.  The Python kernel in the archive computes these chains with exact \code{Fraction} arithmetic.

\subsection{Robust overlapping branch}

Suppose a represented real $x$ is approximated by a rational $r$ with
$\abs{x-r}\le d/4$, where $d>0$, and a downstream proof has two valid local rules:
\[
  x\le q+3d/4\Rightarrow A,
  \qquad
  x\ge q+d/4\Rightarrow B.
\]
Compare $r$ with $q+d/2$.

\begin{itemize}
\item If $r\le q+d/2$, then $x\le r+d/4\le q+3d/4$, so rule $A$ applies.
\item If $r>q+d/2$, then $x\ge r-d/4>q+d/4$, so rule $B$ applies.
\end{itemize}

The true-input ranges overlap.  No decision of $x\le q$ versus $x>q$ is made.  If both branches produce the same observation on the overlap, branch-invariant elimination gives a single theorem.  If they produce values within $\delta$, stable-output elimination gives an error-$\delta$ theorem.

With an additional gap promise
\[
  x\le q-d\quad\text{or}\quad x\ge q+d,
\]
an approximation within $d/3$ decides the exact side.  This is the correct constructive location of exact branching: the gap is part of the input contract.

\subsection{A monotone metastability bound}

Let $(x_n)$ be a nondecreasing rational sequence with $0\le x_n\le M$.  Fix $\eps>0$ and a challenge $g$.  Put
\[
 n_0=0,\qquad n_{t+1}=n_t+g(n_t)+1.
\]
If the window $[n_t,n_t+g(n_t)]$ is not $\eps$-stable, monotonicity gives
\[
 x_{n_{t+1}}-x_{n_t}>\eps.
\]
After $R=\lfloor M/\eps\rfloor+1$ failures, telescoping would give
$x_{n_R}>M$, contradiction.  Thus one of the first $R$ windows is stable.

A quantifier-free invariant is
\[
 \bigl(\text{the first $t$ windows all fail}\bigr)
 \to x_{n_t}>t\eps,
\]
where the finite conjunction of failures is a Boolean fold.  This proves a primitive-recursive stable-window selector.  It does not prove that all windows after the selected one are stable.

\subsection{Finite-net residual minimization}

On $[0,2]$, consider $R(y)=\abs{y^2-2}$.  On this interval,
\[
 \abs{R(y)-R(z)}\le\abs{y^2-z^2}
 \le4\abs{y-z}.
\]
Use the dyadic grid of mesh $2^{-n}$ and choose the least point minimizing the exact rational residual.  Since some grid point is within $2^{-n}$ of $\sqrt2$ (proved internally by least crossing), its residual is at most $4\cdot2^{-n}$.  Hence the argmin has residual no larger.  This is an instance of \Cref{thm:approx-witness}.  Strong monotonicity of $y^2$ on a separated positive interval can additionally turn small residual into a small distance to the positive root.

\subsection{Height-bounded enumeration}

Consider the integer equation
\[
 x^2+y^2=5.
\]
A supplied bound $\max(\abs{x},\abs{y})\le3$ reduces completeness to the finite box $[-3,3]^2$.  A primitive-recursive enumerator tests the equation and returns
\[
 (\pm1,\pm2),\qquad(\pm2,\pm1).
\]
The checker verifies each point, duplicate-freedom under a chosen canonical order, and the completeness implication from the height bound.  The example is elementary, but the architecture is exactly the same for a Runge or Baker height certificate: the hard theorem supplies $B$; the generic bounded filter supplies the complete list.

\subsection{Branch-invariant root data}

For $f(T)=T^2-2$, the two roots cannot be globally ordered by a continuous sign-free operation if coefficients are allowed to pass through collisions in larger families.  Yet the power sum
\[
 s_2=z_1^2+z_2^2
\]
and product $z_1z_2$ are permutation invariant.  A finite relational root-cluster certificate may return either labeling.  Vieta's identities prove that every successful transcript yields $s_2=4$ and $z_1z_2=-2$ (for the exact roots), or the corresponding rational error enclosures for approximations.  \Cref{thm:branch-invariant} removes the label.  This illustrates the general rule: ask for the quotient-level object that the theorem actually uses.

\subsection{Reference implementation coverage}

The included Python tests exercise:
\begin{itemize}
\item bounded first-success search;
\item backward viability and replayed transcript soundness;
\item relational square-root chains at several horizons;
\item branch-invariant terminal outputs;
\item rational error-DAG propagation;
\item monotone metastability search;
\item height/list conversion and finite enumeration.
\end{itemize}
The tests are exact finite checks.  They do not prove the universal theorems and are not an iPRA proof checker.

\section{Theorem ledger, limitations, and validation status}
\label{sec:status}

\subsection{Main claims}

\begin{longtable}{@{}p{0.44\textwidth}p{0.48\textwidth}@{}}
\toprule
Claim & Status in this archive \\
\midrule
\endhead
Bounded Booleanization and least bounded witness & Fully developed human proof; object-library path specified \\
Deterministic finite-trace reflection & Fully developed human proof; source-level \code{PRA.qfInd} wrapper supplied; not compiler-checked here \\
Adaptive finite instantiation of name contracts & Fully developed human proof and invariant design \\
Finite-use viability and relational composition & Fully developed human proof; exact Python implementation/tests; Lean semantic blueprint \\
Relational certificate reflection & Fully developed human proof; checker format specified \\
Branch-invariant and diameter-stable elimination & Fully developed human proofs \\
Regular-name arithmetic and guarded inverse & Fully developed quantitative arguments; object arithmetic library remains to be built \\
Precision-DAG composition & Fully developed human proof; Python exact implementation \\
Rank-to-metastability & Fully developed human proof; Python finite test \\
Root clusters and finite-use consumers & Generic theorem and certificate interface proved; domain-specific root-count producers remain inputs \\
Exact Puiseux branch from arbitrary Cauchy coefficients & Not claimed; retained as obstructed \\
Finite-net transfer and approximate witness search & Fully developed human proofs \\
Nested balls produce a regular name & Fully developed human proof \\
Height bound $\leftrightarrow$ complete finite list & Fully developed human proof; exact finite implementation \\
General Faltings/Siegel/Mordell--Weil producer & Not claimed; isolated as missing height/covering/generator theorem \\
Full \code{Prf PRAAx} relational library & Not completed or compiled \\
Native export of new theorems & Not claimed before object derivations exist \\
\bottomrule
\end{longtable}

\subsection{What has genuinely been resolved}

The report resolves a structural gap between prior deterministic and relational frameworks.  It identifies the exact class of downstream proofs that can consume multivalued evidence without a global selector: bounded adaptive consumers with a finite viability certificate.  It proves that class closed under composition and shows two ways to erase the internal branch---exact observational confluence and quantitative diameter stability.  This directly repairs many formulations involving approximate roots, projective charts, local pivots, and finite covers.

It also gives a more concrete generic object-language wrapper than the existing plans: once theorem-specific base, step, and finish derivations are provided in normalized form, the use of \code{PRA.qfInd}, specialization at the computed bound, cut, and universal introduction is uniform.  The supplied file is not claimed compiled, so this is an architectural and source-level advance rather than a verified repository theorem.

\subsection{What has not been resolved}

The following boundaries remain.

\begin{enumerate}[label=\arabic*.]
\item An infinite coherent branch is not derived from finite viability at every horizon.  A contraction, nested isolation, exact splitting, or uniqueness modulus must supply coherence.
\item Qualitative nonzero does not produce a rational apartness bound.  The inverse theorem is conditional on such a bound.
\item A metastability bound does not become an eventual convergence rate without a tail certificate.
\item Finite root clusters do not become individually named exact Puiseux coefficients for arbitrary Cauchy coefficients.
\item A finite verifier does not produce a missing height bound, zero-free region, covering list, or Mordell--Weil basis.
\item The Lean source has not been built in this artifact environment.  Any declaration whose proof is represented only as a goal proposition or comment remains unfinished.
\end{enumerate}

\subsection{Trust boundary}

The intended trusted chain is:
\[
 \text{small object-language checker soundness}
 \;\Longrightarrow\;
 \Prf\ \PRAAx\ \Gamma\ \varphi
 \;\Longrightarrow\;
 \NPrf\ (\text{translated sequent}).
\]
External certificate generators, numerical optimizers, interval solvers, and precision schedulers may be untrusted.  Their outputs are accepted only after the finite checker verifies every field used by the soundness theorem.

A JSON schema or Lean structure describing a theorem packet is not part of the logical trusted base unless a theorem proves that accepted records imply the object formula.

\subsection{Artifact validation}

The archive records:
\begin{itemize}
\item the exact Python unit-test output;
\item PDF build and text-preflight logs;
\item rendered-page inspection notes;
\item SHA-256 hashes of delivered files;
\item a source audit identifying the repository declarations and packets used as prior art;
\item a machine-readable obstacle matrix and theorem ledger.
\end{itemize}
The Python implementation uses only the standard library and exact integers/\code{Fraction}.  Passing tests checks implementation behavior on finite instances, not the universal human proofs.

\section{Further research questions}
\label{sec:research}

\begin{enumerate}[label=\arabic*.]
\item \textbf{Compile the generic trace wrapper.}  Resolve the remaining de Bruijn substitution normalizations in \code{ObjectTraceReflection.lean}, then prove a constructor from raw (D0)--(D2) local obligations.

\item \textbf{Object-level relational viability.}  Implement rectangular and sparse-certificate versions of the backward viability table as \code{Prf PRAAx} derivations.  Compare proof size and checker complexity.

\item \textbf{A finite-use effect system.}  Give a typed syntax of deterministic and relational certificate transducers with a verified compiler to one trace system.  Prove associativity and observational equivalence inside Lean.

\item \textbf{Automatic finite-query extraction.}  Instrument a semantic evaluator so that it returns both its answer and the finite set of name indices used.  Verify a compiler that turns this trace into the object-language good-table invariant.

\item \textbf{Canonical arithmetic kernel.}  Select one repository implementation of pairs, lists, signed integers, rationals, dyadics, and bounded folds.  Prove a stable API so proof packets stop rebuilding incompatible encodings.

\item \textbf{Precision-DAG certificates.}  Design a compact certificate format for local Lipschitz bounds, magnitude guards, and error budgets.  Verify a scheduler checker rather than trusting generated precision formulas.

\item \textbf{Root-cluster library.}  Formalize real Sturm clusters first, then complex disk counts, multiplicity, and robust coefficient perturbations.  Make unordered/symmetric elimination a standard theorem.

\item \textbf{Projective and spectral data.}  Formalize projective points, subspaces, eigenspace projectors, and separated spectral clusters.  Quantify exactly which Euler--Rodrigues and linear-algebra packets become unconditional at the quotient level.

\item \textbf{Certified compact domains.}  Build reusable nets and coverage proofs for boxes, simplices, spheres, projective spaces, compact matrix groups, and finite probability simplices.

\item \textbf{Rank synthesis.}  Extract bounded descent ranks automatically from energy-decrement, exceptional-set, and finite-variation arguments, while preventing accidental claims of a final convergence rate.

\item \textbf{Analytic certificate taxonomy.}  Identify finite formats for zero-free regions, contour lower bounds, saddle regions, exceptional-zero alternatives, and explicit-formula remainder tables.

\item \textbf{Diophantine producer taxonomy.}  Classify Runge, Baker, descent, Chabauty, Mordell--Weil sieve, and local-cover methods by their produced certificate and generic checker.

\item \textbf{Obstruction library.}  Formalize reusable reductions showing that certain exact selectors imply LLPO, LPO, Markov principles, or a halting decision.  Use them as a pre-formalization linter for theorem statements.

\item \textbf{End-to-end pilot.}  Complete one substantial packet with all layers: fixed evaluators, finite relational certificates, object-level \code{Prf PRAAx}, native export, axiom audit, and independent exact tests.  A separated-real Euler--Rodrigues theorem or a finite-law Berry--Esseen core would exercise most modules.

\item \textbf{Proof-size study.}  Compare ordinary \code{PRAAx}, the native scalar core, public guarded certificates, and a specialized finite-use checker by proof-term size, checker source size, and compression.

\item \textbf{Repository-wide status automation.}  Generate a dashboard from packet metadata that distinguishes human proof, executable finite evidence, compiled host theorem, embedded object derivation, native export, and unresolved producer.
\end{enumerate}

\appendix
\section{Normal forms for theorem statements}
\label{app:normal-forms}

The following normal forms make an intended iPRA theorem auditable.

\subsection{Deterministic producer}

A deterministic construction should be stated as
\[
 \forall x\;\bigl(\mathsf{Pre}(x)=1\to
 \mathsf{Post}(x,F(x))=1\bigr),
\]
where $F$, \code{Pre}, and \code{Post} are fixed primitive-recursive functions or finite-observation predicates.  The proof artifact identifies the trace, invariant, and local obligations.

\subsection{Precision-indexed represented output}

A represented construction should be stated
\[
 \forall p\;\bigl(H(p)\to\forall k\;
 \mathsf{Obs}(p,k,A(p,k))=1\bigr),
\]
with a primitive-recursive query planner for each free $k$.  The proof order is $p,k$ free; finite trace; output observation; universal introduction over $k$.

\subsection{Relational output}

A multivalued producer should expose
\[
 \forall x,h\;\bigl(\mathsf{Evidence}(x,h)=1\to
 \Viable(x,h,\mathsf{start}(x))=1\bigr).
\]
A consumer theorem then concludes a successful transcript, an invariant observation, or a diameter-stable output.  It should not silently replace viability by an infinite selector.

\subsection{Conditional certificate theorem}

For an externally generated certificate $c$:
\[
 \forall x,c\;\bigl(\mathsf{Check}(x,c)=1\to
 \mathsf{Conclusion}(x,\mathsf{output}(c))\bigr).
\]
The producer status is recorded separately.  A theorem that merely defines \code{Check} is not the soundness theorem.

\subsection{Height/list theorem}

A finite Diophantine result should name the producer:
\[
 \mathsf{CurveHyp}(x)\to
 \exists B\le H(x)\;\forall z\,(P(x,z)\to\mathsf{height}(z)\le B),
\]
or, constructively stronger and more useful, give a concrete $B(x)$.  The generic enumerator theorem begins only after this field is available.

\subsection{Exact choice under a gap}

An exact branch theorem should have the form
\[
 \mathsf{Gap}(x,m)=1\to
 \mathsf{ExactBranch}(x,\mathsf{select}(x,m))=1.
\]
Without \code{Gap}, use overlapping branches, a finite relation, or a quotient output.

\section{Repository source audit and novelty boundary}
\label{app:source-audit}

The repository was inspected read-only at main commit
\path{0edaf721b175de6aa219755df18b1588be803a16}.  The most relevant paths were:

\begin{itemize}
\item \path{Logic/PrimitiveRecursiveArithmetic/README.md} and \path{docs/proofs/README.md}: checked/unverified status boundary and packet catalog.
\item \path{ConciseIPRA/Lean/ConciseIPRA/PR.lean}: \code{PR}, \code{PTm}, \code{PRAAx}, \code{PRA.qfInd}, and defining equations.
\item \path{ConciseIPRA/Lean/ConciseIPRA/Deduction.lean}: \code{Prf}, substitution, weakening, cut, equality, and quantifier rules.
\item \path{ConciseIPRA/Lean/ConciseIPRA/NativeBiInterpretation.lean}: \code{native_iff_pra}, \code{pra_iff_native}, and round trips.
\item \path{docs/proofs/foundations/ipra_embedding_framework}: finite observations, simple analytic branches, stability, finite algebra certificates.
\item \path{docs/proofs/foundations/ipra_evidence_framework}: finite relational outputs and backward finite compatibility.
\item \path{docs/proofs/foundations/ipra_finite_observation}: finite observation compiler, degree-uniform factorization, continuity and section bounds.
\item \path{docs/proofs/foundations/ipra_framework}: finite contracts, represented arithmetic, metastability, local branches, and power gaps.
\item \path{docs/proofs/foundations/ipra_framework_2026-11-15}: certificate-and-modulus summary and obstruction discipline.
\item \path{docs/proofs/foundations/ipra_uniform_certificates}: finite-net lifting, separated roots, and precision composition.
\item \path{docs/proofs/foundations/ipra_embedding_framework_2026-11-15}: Certified Finitary Presentations, deterministic finite-trace theorem, robust branches, nets, barriers, and object-language plan.
\item \path{docs/proofs/algebra/puiseux_ipra/evidence/scope_obstruction.md}: exact Puiseux selection obstruction for arbitrary Cauchy coefficients.
\item \path{docs/proofs/diophantine-number-theory/faltings_ipra} and \path{siegel_ipra}: height/list boundaries and effective subcases.
\end{itemize}

The report credits and consolidates the prior techniques.  The principal additional contribution is the consumer-oriented finite-use relational composition theorem, including adaptive state, certificate reflection, branch-invariant elimination, and quantitative diameter elimination.  The source-level object trace wrapper is also more explicit than the prior plans, but is not claimed compiled.

No priority claim is made for general bounded dynamic programming, proof mining, metastability, constructive analysis, or finite certificate checking.  The novelty asserted is only the specific synthesis and interface tailored to the inspected ProveIt iPRA program.

\section{Notation and implementation conventions}
\label{app:notation}

\begin{longtable}{@{}p{0.25\textwidth}p{0.67\textwidth}@{}}
\toprule
Symbol & Meaning \\
\midrule
\endhead
$\PR$ & Primitive-recursive, always for a fixed definition unless a bounded program table is explicitly supplied \\
$\B$ & Numerical Booleans $\{0,1\}$ \\
$u,x,p$ & Paired numerical input/certificate/representation codes \\
$k$ & Requested output precision \\
$h,r$ & Interaction horizon or remaining resource \\
$s$ & Encoded deterministic or consumer state \\
$\Viable(x,r,s)$ & Backward Boolean saying a successful finite continuation exists \\
$\Local(x,r,s,a)$ & Local admissibility of answer $a$ \\
$E$ & Observational equivalence, not necessarily literal equality \\
$\delta$ & Diameter/error bound for stable relational outputs \\
$\Prf\ \PRAAx\ \Gamma\ \varphi$ & Object-level derivability in ordinary iPRA \\
$\NPrf$ & Derivability in the native scalar calculus \\
\bottomrule
\end{longtable}

All finite sequences and graphs may be encoded by any standard primitive-recursive pairing/list scheme with explicit bounds.  Proofs use only closure properties and bounded lookup; no specific optimal coding is mathematically essential.  Implementations should nevertheless share one canonical repository library to avoid repeated conversion proofs.

Rational inequalities are intended to be reduced to signed-integer cross products with positive denominators.  Every use of an order decision is therefore a decision about exact integers, not an arbitrary represented real.


\begin{thebibliography}{99}
\bibitem{BridgesRichman}
D. Bridges and F. Richman, \emph{Varieties of Constructive Mathematics}, London Mathematical Society Lecture Note Series 97, Cambridge University Press, 1987.

\bibitem{Kohlenbach}
U. Kohlenbach, \emph{Applied Proof Theory: Proof Interpretations and Their Use in Mathematics}, Springer Monographs in Mathematics, Springer, 2008.

\bibitem{KohlenbachLeustean}
U. Kohlenbach and L. Leu\c{s}tean, Effective metastability of Halpern iterates in CAT(0) spaces, \emph{Advances in Mathematics} 231 (2012), 2526--2556.

\bibitem{MannaaCoquand}
B. Mannaa and T. Coquand, Dynamic Newton--Puiseux theorem, \emph{Journal of Logic and Analysis} 5 (2013), Paper 5.

\bibitem{ProveIt}
V. Reshetnikov, \emph{ProveIt}, repository directory \code{Logic/PrimitiveRecursiveArithmetic}, inspected at main commit \code{0edaf721b175de6aa219755df18b1588be803a16}, 15 November 2026 (Pacific time).
\end{thebibliography}

\end{document}
